\documentclass[12pt, a4paper]{article}
\usepackage[T2A]{fontenc}
\usepackage[utf8]{inputenc}
\usepackage[russian]{babel}
\usepackage{amsmath, amssymb, amsthm}
\usepackage{geometry}
\usepackage{hyperref}
\usepackage{enumitem}

\geometry{left=2.5cm, right=2.5cm, top=2.5cm, bottom=2.5cm}

\newtheorem{theorem}{Теорема}
\newtheorem{lemma}[theorem]{Лемма}
\newtheorem{corollary}[theorem]{Следствие}
\theoremstyle{definition}
\newtheorem{definition}{Определение}
\newtheorem{example}{Пример}
\theoremstyle{remark}
\newtheorem*{remark}{Замечание}

\newcommand{\Z}{\mathbb{Z}}
\newcommand{\Q}{\mathbb{Q}}
\newcommand{\R}{\mathbb{R}}
\renewcommand{\C}{\mathbb{C}}
\newcommand{\F}{\mathbb{F}}
\newcommand{\Ker}{\mathrm{Ker}\,}
\renewcommand{\Im}{\mathrm{Im}\\,}

\title{\textbf{Ответы на вопросы к экзамену по ТФКП}\\\large 1-24, Весна 2026}
\author{}
\date{}

\begin{document}
	\maketitle
	\tableofcontents
	\newpage
	
%──────────────────────────────────────────────────────────────────────────────
\section{Непрерывность и дифференцируемость функций комплексного переменного, их связь. Теорема Коши–Римана. Голоморфные функции}

Определение 2.1. Функция $f : U(z_0) \to \mathbb{C}$ называется дифференцируемой в точке $z_0$, если существует $A \in \mathbb{C}$ такое, что
\begin{equation}
	f(z) - f(z_0) = A(z - z_0) + o(z - z_0), \quad z \to z_0.
\end{equation}
Линейная форма $df(z_0) : z \mapsto Az$ называется дифференциалом функции $f$ в точке $z_0$.

Определение 2.2. Производной функции $f : U(z_0) \to \mathbb{C}$ в точке $z_0$ называется предел (если он существует)
\begin{equation}
	f'(z_0) = \lim_{z \to z_0} \frac{f(z) - f(z_0)}{z - z_0}.
\end{equation}

Как и в действительном случае, справедливо
\textbf{Предложение 2.3.} Функция $f : U(z_0) \to \mathbb{C}$ дифференцируема в $z_0$ тогда и только тогда, когда она имеет в $z_0$ производную. При этом в (2.1) имеем $A = f'(z_0)$.

Пусть функция $f(z)$ определена в некоторой окрестности точки $z_0 \in \mathbb{C}$.

\textbf{Определение 2.3.} Функция $f(z)$ называется \textbf{непрерывной в точке} $z_0$, если предел функции в этой точке существует и равен значению функции в этой точке:
\[
\lim_{z \to z_0} f(z) = f(z_0).
\]
На языке $\varepsilon-\delta$ это означает, что для любого $\varepsilon > 0$ существует такое $\delta > 0$, что для всех $z$, удовлетворяющих условию $|z - z_0| < \delta$, выполняется неравенство:
\[
|f(z) - f(z_0)| < \varepsilon.
\]

\textbf{Теорема 2.5.} Если функция $f(z)$ является $\mathbb{C}$-дифференцируемой в точке $z_0$, то она непрерывна в этой точке.

\textbf{Доказательство.} Пусть функция $f(z)$ дифференцируема в точке $z_0$. По определению комплексной производной, приращение функции $\Delta f = f(z) - f(z_0)$ в окрестности этой точки можно представить в виде:
\[
f(z) - f(z_0) = f'(z_0)(z - z_0) + o(z - z_0), \quad z \to z_0,
\]
где $f'(z_0) \in \mathbb{C}$ --- конечное комплексное число, а $o(z - z_0)$ --- бесконечно малая функция более высокого порядка по сравнению с $(z - z_0)$, то есть $\lim_{z \to z_0} \frac{o(z - z_0)}{z - z_0} = 0$.

Перейдем к пределу при $z \to z_0$ в выражении для приращения функции:
\[
\lim_{z \to z_0} \left( f(z) - f(z_0) \right) = \lim_{z \to z_0} \left( f'(z_0)(z - z_0) + o(z - z_0) \right).
\]
Используя линейные свойства предела, получаем:
\[
\lim_{z \to z_0} \left( f(z) - f(z_0) \right) = f'(z_0) \cdot \lim_{z \to z_0} (z - z_0) + \lim_{z \to z_0} o(z - z_0).
\]
Так как $\lim_{z \to z_0} (z - z_0) = 0$ и $\lim_{z \to z_0} o(z - z_0) = 0$, правая часть равенства обращается в ноль:
\[
\lim_{z \to z_0} \left( f(z) - f(z_0) \right) = f'(z_0) \cdot 0 + 0 = 0.
\]
Отсюда следует, что:
\[
\lim_{z \to z_0} f(z) = f(z_0),
\]
что означает непрерывность функции $f(z)$ в точке $z_0$. Теорема доказана. $\quad \blacksquare$

\textbf{Замечание.} Обратное утверждение неверно: непрерывность является лишь необходимым, но не достаточным условием дифференцируемости.

\subsection{Условия Коши --- Римана}

\textbf{Теорема 2.4.} Функция $f = (x + iy \mapsto u(x,y) + iv(x,y)) \in \mathbb{C}^{U(z_0)}$ является $\mathbb{C}$-дифференцируемой в $z_0 = x_0 + iy_0$ тогда и только тогда, когда функция $F = ((x,y) \mapsto (u(x,y), v(x,y))) \in (\mathbb{R}^2)^{U(x_0,y_0)}$ дифференцируема в $(x_0,y_0)$ и выполнены равенства
\[
\frac{\partial u(x_0, y_0)}{\partial x} = \frac{\partial v(x_0, y_0)}{\partial y}, \quad \frac{\partial u(x_0, y_0)}{\partial y} = -\frac{\partial v(x_0, y_0)}{\partial x}.
\]
В этом случае
\[
f'(z_0) = \frac{\partial u(x_0, y_0)}{\partial x} + i \frac{\partial v(x_0, y_0)}{\partial x}.
\]

\textbf{Доказательство.} 

$(\Leftarrow)$ Пусть функция $F = (u, v)$ дифференцируема в точке $(x_0, y_0)$ как функция двух вещественных переменных и выполнены условия Коши–Римана. 

Из $\mathbb{R}^2$-дифференцируемости приращения функций $u$ и $v$ в окрестности точки $(x_0, y_0)$ можно представить в виде:
\[
\Delta u = \frac{\partial u}{\partial x}\Delta x + \frac{\partial u}{\partial y}\Delta y + o(\rho),
\]
\[
\Delta v = \frac{\partial v}{\partial x}\Delta x + \frac{\partial v}{\partial y}\Delta y + o(\rho),
\]
где $\rho = \sqrt{\Delta x^2 + \Delta y^2} = |\Delta z|$, а все частные производные берутся в точке $(x_0, y_0)$.

Составим комплексное приращение функции $\Delta f = \Delta u + i \Delta v$:
\[
\Delta f = \left( \frac{\partial u}{\partial x}\Delta x + \frac{\partial u}{\partial y}\Delta y \right) + i \left( \frac{\partial v}{\partial x}\Delta x + \frac{\partial v}{\partial y}\Delta y \right) + o(|\Delta z|).
\]
Используя условия Коши–Римана, заменим $\frac{\partial u}{\partial y}$ на $-\frac{\partial v}{\partial x}$, а $\frac{\partial v}{\partial y}$ на $\frac{\partial u}{\partial x}$:
\[
\Delta f = \left( \frac{\partial u}{\partial x}\Delta x - \frac{\partial v}{\partial x}\Delta y \right) + i \left( \frac{\partial v}{\partial x}\Delta x + \frac{\partial u}{\partial x}\Delta y \right) + o(|\Delta z|).
\]
Сгруппируем слагаемые при частных производных:
\[
\Delta f = \frac{\partial u}{\partial x} (\Delta x + i\Delta y) + i \frac{\partial v}{\partial x} (\Delta x + i\Delta y) + o(|\Delta z|).
\]
Вынося общую комплексную скобку $\Delta z = \Delta x + i\Delta y$, получаем:
\[
\Delta f = \left( \frac{\partial u}{\partial x} + i \frac{\partial v}{\partial x} \right) \Delta z + o(|\Delta z|).
\]
Обозначив комплексное число в скобках как $A = \frac{\partial u}{\partial x} + i \frac{\partial v}{\partial x} \in \mathbb{C}$, мы приходим к равенству:
\[
f(z) - f(z_0) = A(z - z_0) + o(z - z_0), \quad z \to z_0.
\]
Согласно Определению 2.1, это означает, что функция $f$ является $\mathbb{C}$-дифференцируемой в точке $z_0$, и $f'(z_0) = A$. Достаточность доказана. $\quad \blacksquare$

$(\Rightarrow)$ Пусть функция $f$ является $\mathbb{C}$-дифференцируемой в точке $z_0 = x_0 + iy_0$. Это означает, что существует комплексное число $A = A_1 + iA_2 \in \mathbb{C}$ (где $A_1, A_2 \in \mathbb{R}$), такое что приращение функции $\Delta f$ в окрестности $z_0$ представимо в виде:
\[
f(z_0 + \Delta z) - f(z_0) = \Delta f = A \Delta z + o(|\Delta z|), \quad \Delta z \to 0.
\]
Расписав комплексные переменные через их вещественные и мнимые части $\Delta z = \Delta x + i\Delta y$ и $\Delta f = \Delta u + i \Delta v$, подставим их в равенство дифференцируемости:
\[
\Delta u + i \Delta v = (A_1 + iA_2)(\Delta x + i\Delta y) + o(|\Delta z|).
\]
Выполним комплексное умножение в правой части:
\[
\Delta u + i \Delta v = (A_1 \Delta x - A_2 \Delta y) + i(A_2 \Delta x + A_1 \Delta y) + o(|\Delta z|).
\]
Выделяя вещественную и мнимую части в полученном равенстве, получаем систему для приращений функций $u$ и $v$:
\[
\begin{cases}
	\Delta u = A_1 \Delta x - A_2 \Delta y + o(|\Delta z|), \\
	\Delta v = A_2 \Delta x + A_1 \Delta y + o(|\Delta z|).
\end{cases}
\]
Поскольку $|\Delta z| = \sqrt{\Delta x^2 + \Delta y^2} = \rho$, то $o(|\Delta z|) = o(\rho)$. Линейность главной части приращений относительно $\Delta x$ и $\Delta y$ гарантирует, что вещественные функции $u(x,y) и v(x,y)$ дифференцируемы в точке $(x_0, y_0)$ как функции двух вещественных переменных.

Из единственности коэффициентов линейной формы вещественного дифференциала (которые равны соответствующим частным производным) немедленно следуют равенства:
\[
\begin{cases}
	\dfrac{\partial u}{\partial x} = A_1, \quad \dfrac{\partial u}{\partial y} = -A_2, \\
	\dfrac{\partial v}{\partial x} = A_2, \quad \dfrac{\partial v}{\partial y} = A_1.
\end{cases}
\]
Сопоставляя коэффициенты через $A_1$ и $A_2$, получаем условия Коши — Римана:
\[
\frac{\partial u}{\partial x} = \frac{\partial v}{\partial y} = A_1, \quad \frac{\partial u}{\partial y} = -\frac{\partial v}{\partial x} = -A_2.
\]
Наконец, выражая производную $f'(z_0) = A = A_1 + iA_2$ через найденные частные производные, находим:
\[
f'(z_0) = \frac{\partial u(x_0, y_0)}{\partial x} + i \frac{\partial v(x_0, y_0)}{\partial x}.
\]
Необходимость доказана. $\quad \blacksquare$

Определение 2.5. Функция $f$ называется голоморфной, или аналитической в точке $z_0 \in \mathbb{C}$, если она дифференцируема в некоторой её окрестности. Мы говорим, что функция $f$ голоморфна (аналитична) на множестве $M$, если существует её голоморфное продолжение на некоторое открытое множество $D \supset M$.

%──────────────────────────────────────────────────────────────────────────────
\newpage
\section{Геометрический смысл комплексной производной. Конформные отображения, связь конформности и дифференцируемости, примеры}
\textbf{Предложение 2.6.} Любая голоморфная в точке $a \in \mathbb{C}$ функция $f$ определяет линейное отображение касательных векторов $\eta = f'(a)\xi$, где $\eta$ --- касательный вектор в точке $a$, $\xi$ --- касательный вектор в точке $f(a)$. Это отображение представляет собой растяжение с коэффициентом $|f'(a)|$ и углом поворота $\arg f'(a)$.

\textbf{Доказательство Предложения 2.6.} 
Пусть через точку $a$ в $Z$-плоскости проходит гладкая кривая $\gamma: z(t) = x(t) + iy(t)$, $t \in [\alpha, \beta]$. Пусть $z(t_0) = a$. Касательный вектор к кривой $\gamma$ в точке $a$ задается комплексным числом 
\[
\xi = z'(t_0) = \dot{x}(t_0) + i\dot{y}(t_0).
\]
Его аргумент $\arg\xi = \varphi$ задает угол наклона касательной к вещественной оси.

При отображении $w = f(z)$ образ кривой $\gamma$ есть кривая $\Gamma = f(\gamma)$, задаваемая уравнением $w(t) = f(z(t))$. По правилу дифференцирования сложной функции, касательный вектор к кривой $\Gamma$ в точке $w_0 = f(a)$ равен:
\[
\eta = w'(t_0) = f'(z(t_0)) \cdot z'(t_0) = f'(a)\xi.
\]
Если $f'(a) \neq 0$, то переходя к модулю и аргументу этого равенства, получаем:
\[
|\eta| = |f'(a)| \cdot |\xi|, \quad \arg\eta = \arg f'(a) + \arg\xi.
\]
Это означает, что:
1) Длина любого касательного вектора $\xi$ при отображении умножается на одно и то же число $k = |f'(a)|$ (свойство постоянства растяжений).
2) Направление любого касательного вектора поворачивается на один и тот же угол $\alpha = \arg f'(a)$ (свойство сохранения углов). 
Линейное отображение $\xi \mapsto f'(a)\xi$ является композицией поворотной гомотетии и параллельного переноса. $\quad \blacksquare$

\vspace{1em}
\subsection{Определение конформности}

В геометрии и теории ФКП принято более развернутое определение, инвариантное к выбору системы координат.

\textbf{Определение 2.6' (Геометрическое).} Отображение окрестности точки $z_0$ называется \textit{конформным} (или сохраняющим форму), если оно обладает свойствами:
\begin{enumerate}
	\item \textbf{Сохранения углов:} угол между любыми двумя непрерывными гладкими кривыми $\gamma_1$ и $\gamma_2$, пересекающимися в точке $z_0$, равен по величине и направлению (ориентации) углу между их образами $\Gamma_1$ и $\Gamma_2$ в точке $w_0 = f(z_0)$.
	\item \textbf{Постоянства растяжений:} отношение длины бесконечно малого вектора в окрестности точки $w_0$ к длине его прообраза в окрестности $z_0$ стремится к пределу, зависящему только от точки $z_0$, но не от направления вектора.
\end{enumerate}

\vspace{0.5em}
\textbf{Следствие 2.7 (Теорема о конформности голоморфного отображения).} 
Пусть функция $f(z)$ является голоморфной в открытой области $U \subset \mathbb{C}$ ($f \in C^{\omega}(U)$). Тогда:
\begin{enumerate}
	\item Отображение, осуществляемое функцией $f(z)$, является \textit{конформным в точке} $z_0 \in U$ тогда и только тогда, когда $f'(z_0) \neq 0$.
	\item Если условие $f'(z) \neq 0$ выполняется во всех точках области $U$, то отображение $f: U \to \mathbb{C}$ является \textit{локально конформным} в этой области.
	\item Если, кроме того, функция $f(z)$ является \textit{однолистной} (взаимно однозначной) в области $U$, то она осуществляет \textit{конформное отображение области $U$ на область $V = f(U)$}.
\end{enumerate}
\textbf{Доказательство Следствия 2.7 (Связь конформности и голоморфности).} 

\textit{Необходимость.} Пусть отображение $f$ конформно в точке $z_0$. Из постоянства растяжений следует, что при $\Delta z \to 0$ модуль отношения $|\Delta w| / |\Delta z|$ стремится к конечному пределу $k$, независимо от того, по какому направлению $\Delta z \to 0$. Из сохранения углов следует, что аргумент $\arg(\Delta w / \Delta z)$ также стремится к фиксированному углу $\alpha$. 

Следовательно, существует независимый от направления предел комплексного отношения:
\[
\lim_{\Delta z \to 0} \frac{\Delta w}{\Delta z} = k e^{i\alpha} = f'(z_0).
\]
Это доказывает $\mathbb{C}$-дифференцируемость (голоморфность) функции $f(z)$ в точке $z_0$. Так как по условию свойства конформности выполняются, масштаб $k = |f'(z_0)| \neq 0$, то $f'(z_0) \neq 0$.

\textit{Достаточность.} Пусть $f(z)$ голоморфна в области $U$ ($f \in C^{\omega}(U)$) и $f'(z) \neq 0$. Из Предложения 2.6 следует, что для любых двух кривых $\gamma_1$ и $\gamma_2$, пересекающихся под углом $\Delta \varphi = \arg \xi_2 - \arg \xi_1$, их образы будут пересекаться под углом:
\[
\Delta \Phi = \arg \eta_2 - \arg \eta_1 = (\arg f'(z_0) + \arg \xi_2) - (\arg f'(z_0) + \arg \xi_1) = \arg \xi_2 - \arg \xi_1 = \Delta \varphi.
\]
Угол сохраняется как по величине, так и по направлению отсчета. Постоянство растяжений также обеспечено равенством $|\eta|/|\xi| = |f'(z_0)| = \text{const}$ для всех направлений в данной точке. Отображение конформно. $\quad \blacksquare$

%──────────────────────────────────────────────────────────────────────────────
\newpage
\section{Дробно–линейные функции, их конформность и групповое свойство}
Определение 3.3. Дробно-линейным отображением (ДЛО) называется функция вида
\begin{equation}
	w : z \mapsto \frac{az + b}{cz + d},
\end{equation}
где $a, b, c, d \in \mathbb{C}$, $ad \neq bc$. При этом $w(\infty) \stackrel{\text{def}}{=} a/c$ и $w(-d/c) \stackrel{\text{def}}{=} \infty$.

\textbf{Предложение 3.5.} Отображение $w$ конформно в каждой точке $\hat{\mathbb{C}}$.

\textbf{Доказательство.} Для неособых точек теорема следует из существования производной
\[
w'(z) = \frac{ad - bc}{(cz + d)^2} \neq 0.
\]
Для доказательства утверждения в особых точках $-d/c$ и $\infty$ достаточно рассмотреть соответственно отображения $1/w$ и $w(1/z)$. $\blacksquare$

\vspace{1em}
\subsection{Групповое свойство дробно-линейных отображений}



\textbf{Теорема 3.6 (О групповом свойстве).} Совокупность всех дробно-линейных отображений образует группу относительно операции суперпозиции (композиции) функций.

\textbf{Доказательство.} Обозначим множество всех ДЛО через $\mathcal{M}$. Для того чтобы $\mathcal{M}$ являлось группой, должны выполняться следующие четыре аксиомы:

1. \textbf{Замкнутость.} Пусть даны два произвольных дробно-линейных отображения:
\[
w_1(z) = \frac{a_1 z + b_1}{c_1 z + d_1} \quad (a_1 d_1 - b_1 c_1 \neq 0), \quad w_2(z) = \frac{a_2 z + b_2}{c_2 z + d_2} \quad (a_2 d_2 - b_2 c_2 \neq 0).
\]
Рассмотрим их композицию $w(z) = w_2(w_1(z))$:
\[
w(z) = \frac{a_2 \left(\dfrac{a_1 z + b_1}{c_1 z + d_1}\right) + b_2}{c_2 \left(\dfrac{a_1 z + b_1}{c_1 z + d_1}\right) + d_2} = \frac{(a_2 a_1 + b_2 c_1)z + (a_2 b_1 + b_2 d_1)}{(c_2 a_1 + d_2 c_1)z + (c_2 b_1 + d_2 d_1)}.
\]
Положим $a = a_2 a_1 + b_2 c_1$, $b = a_2 b_1 + b_2 d_1$, $c = c_2 a_1 + d_2 c_1$, $d = c_2 b_1 + d_2 d_1$. Прямым вычислением проверяется определитель полученного отображения:
\[
ad - bc = (a_2 d_2 - b_2 c_2)(a_1 d_1 - b_1 c_1).
\]
Так как по условию оба сомножителя в правой части отличны от нуля, то $ad - bc \neq 0$. Следовательно, композиция $w_2 \circ w_1$ также является дробно-линейным отображением и принадлежит множеству $\mathcal{M}$.

2. \textbf{Ассоциативность.} Для любых трех отображений $w_1, w_2, w_3 \in \mathcal{M}$ в силу общих свойств суперпозиции отображений выполняется соотношение:
\[
w_3 \circ (w_2 \circ w_1) = (w_3 \circ w_2) \circ w_1.
\]

3. \textbf{Наличие нейтрального элемента.} Роль единицы группы играет тождественное отображение $e(z) = z$, которое получается из общего вида ДЛО при $a=d=1$, $b=c=0$. Очевидно, что $ad - bc = 1 \cdot 1 - 0 \cdot 0 = 1 \neq 0$, то есть $e(z) \in \mathcal{M}$. При этом для любого $w \in \mathcal{M}$ имеем $w \circ e = e \circ w = w$.

4. \textbf{Наличие обратного элемента.} Каждое дробно-линейное отображение $w = \frac{az + b}{cz + d}$ взаимно однозначно отображает $\hat{\mathbb{C}}$ на себя. Разрешая уравнение относительно $z$, находим обратное отображение:
\[
w(cz + d) = az + b \implies z(cw - a) = -dw + b \implies z = w^{-1}(w) = \frac{-dw + b}{cw - a}.
\]
Определитель обратного отображения равен $(-d)(-a) - b \cdot c = ad - bc \neq 0$, следовательно, $w^{-1} \in \mathcal{M}$. При этом $w \circ w^{-1} = w^{-1} \circ w = e$.

Таким образом, множество $\mathcal{M}$ образует группу относительно операции суперпозиции. $\quad \blacksquare$

%──────────────────────────────────────────────────────────────────────────────
\newpage
\section{Дробно–линейные функции, их геометрическая интерпретация и свойство трех точек}
\textbf{Теорема 4.1.} Для любых троек $(z_1, z_2, z_3), (w_1, w_2, w_3) \in \hat{\mathbb{C}}^3$, элементы каждой из которых попарно различны, существует единственное ДЛО $L$ такое, что $L(z_i) = w_i$, $i = 1, 2, 3$.

\textbf{Доказательство.} Доказательство существования начнём с частного случая $(w_1, w_2, w_3) = (0, \infty, 1)$. Тогда
\[
L : z \mapsto \frac{z - z_1}{z - z_2} \cdot \frac{z_3 - z_2}{z_3 - z_1}
\]
является искомым отображением. В общем случае построим два отображения
\[
L_1 : (z_1, z_2, z_3) \mapsto (0, \infty, 1), \quad L_2 : (w_1, w_2, w_3) \mapsto (0, \infty, 1).
\]
Тогда искомое отображение $L \equiv w(z)$ будет композицией $L_2^{-1} \circ L_1$. Оно определяется выражением
\begin{equation}
	\frac{z - z_1}{z - z_2} \cdot \frac{z_3 - z_2}{z_3 - z_1} = \frac{w - w_1}{w - w_2} \cdot \frac{w_3 - w_2}{w_3 - w_1}.
\end{equation}

Докажем единственность отображения $L$. Пусть
\[
L_1 : z \mapsto \frac{z - z_1}{z - z_2} \cdot \frac{z_3 - z_2}{z_3 - z_1}, \quad L_2 : w \mapsto \frac{w - w_1}{w - w_2} \cdot \frac{w_3 - w_2}{w_3 - w_1},
\]
а $\lambda$ — произвольное ДЛО $(z_1, z_2, z_3) \mapsto (w_1, w_2, w_3)$. Тогда $\mu = L_2 \circ \lambda \circ L_1^{-1}$ — ДЛО с неподвижными точками $0, \infty, 1$. Из условия $\mu(\infty) = \infty$ получаем, что $\mu$ — целая линейная функция, т. е. $\mu(\xi) = a\xi + b$. Подставляя $\mu(0) = 0$, $\mu(1) = 1$, получим, что $a = 1$ и $b = 0$. Таким образом, $\mu = \text{id}_{\hat{\mathbb{C}}}$, что и доказывает единственность. $\quad \blacksquare$

Геометрический смысл ДЛО $\hat{\mathbb{C}} \to \hat{\mathbb{C}}$ полностью описывается двумя фундаментальными свойствами: круговым и свойством сохранения симметрии.

Для их анализа разложим общее отображение $w = \frac{az + b}{cz + d}$ при $c \neq 0$ в композицию элементарных отображений. Выделим целую часть:
\[
w = \frac{\dfrac{a}{c}(cz + d) - \dfrac{ad}{c} + b}{cz + d} = \frac{a}{c} + \frac{bc - ad}{c(cz + d)} = \frac{a}{c} + \frac{bc - ad}{c^2} \cdot \frac{1}{z + \dfrac{d}{c}}.
\]
Отсюда видно, что любое ДЛО является суперпозицией следующих четырех простейших отображений:
\begin{enumerate}
	\item $w_1 = z + \dfrac{d}{c}$ --- параллельный перенос (сдвиг);
	\item $w_2 = \dfrac{1}{w_1}$ --- инверсия (круговая симметрия относительно единичной окружности совмещенная с комплексным сопряжением);
	\item $w_3 = \frac{bc - ad}{c^2} \cdot w_2$ --- поворотная гомотетия (растяжение и поворот);
	\item $w = w_3 + \dfrac{a}{c}$ --- финальный параллельный перенос.
\end{enumerate}
Поскольку отображения 1, 3 и 4 представляют собой линейные преобразования (поворот, сдвиг, масштаб), их геометрические свойства очевидны. Главным и наиболее содержательным элементом ДЛО является отображение инверсии $w = 1/z$.

%──────────────────────────────────────────────────────────────────────────────
\newpage
\section{Дробно–линейные функции, их геометрические свойства: круговое свойство и сохранение симметрии}
\textbf{Теорема 4.2.} ДЛО отображают обобщённые окружности в обобщённые окружности.

\textbf{Доказательство.} Пусть
\[
L : z \mapsto \frac{az + b}{cz + d}.
\]
1) Рассмотрим сначала случай $c = 0$. Тогда $a \neq 0$ и $L = L_2 \circ L_1$, где
\[
L_1(z) = z + \frac{b}{d}, \quad L_2(z) = az.
\]
В силу предложения 2.6 отображение $L$ является композицией сдвига и растяжения c поворотом, а такие отображения сохраняют обобщённые окружности.

2) Пусть $c \neq 0$. Тогда
\[
L(z) = \frac{a}{c} - \frac{ad - bc}{c(cz + d)} \stackrel{\text{def}}{=} A + \frac{B}{z + D}
\]
и $L = L_3 \circ L_2 \circ L_1$, где
\[
L_1(z) = z + D, \quad L_2(z) = \frac{1}{z}, \quad L_3(z) = A + Bz.
\]
Мы докажем наше утверждение, если покажем, что $L_2$ сохраняет обобщённые окружности.

Обобщённые окружности на плоскости $\mathbb{R}^2$ описываются уравнениями
\[
E(x^2 + y^2) + F_1x + F_2y + C = 0,
\]
где $E^2 + F_1^2 + F_2^2 + C^2 \neq 0$. На комплексной плоскости эти уравнения принимают вид
\[
E z\bar{z} + \frac{F_1 - F_2i}{2}z + \frac{F_1 + F_2i}{2}\bar{z} + C \stackrel{\text{def}}{=} E z\bar{z} + Fz + \bar{F}\bar{z} + C = 0.
\]
Если точка $z = 1/w$ удовлетворяет этому уравнению, то
\[
C w\bar{w} + F\bar{w} + \bar{F}w + E = 0.
\]
Это уравнение также задаёт обобщённую окружность. $\quad \blacksquare$

\vspace{1em}
\textbf{Теорема 4.4.} Если $L$ — ДЛО, а точки $z_1$ и $z_2$ симметричны относительно обобщённой окружности $\Gamma$, то $w_1 = L(z_1)$ и $w_2 = L(z_2)$ симметричны относительно $L(\Gamma)$.

\textbf{Доказательство.} Пусть $\gamma \ni w_1, w_2$ — произвольная обобщённая окружность; тогда $L^{-1}(\gamma)$ — обобщённая окружность, проходящая через $z_1$ и $z_2$. Поскольку $L^{-1}(\gamma) \perp \Gamma$, а ДЛО сохраняет углы, имеем $\gamma \perp L(\Gamma)$. $\quad \blacksquare$

%──────────────────────────────────────────────────────────────────────────────
\newpage
\section{Стереографическая проекция. Расширенная комплексная плоскость и её топология. Бесконечно удаленная точка, её окрестности. Угол между кривыми в бесконечности. Дифференцируемость и конформность в бесконечности. Дробно–линейные функции как отображения расширенной комплексной плоскости}

Положим $\hat{\mathbb{C}} \stackrel{\text{def}}{=} \mathbb{C} \cup \{\infty\}$, где $\infty$ --- произвольный элемент, называемый бесконечно удалённой точкой. По определению,
\[
z + \infty \stackrel{\text{def}}{=} \infty, \quad \frac{z}{\infty} \stackrel{\text{def}}{=} 0 \quad (z \neq \infty);
\]
\[
z \cdot \infty \stackrel{\text{def}}{=} \infty, \quad \frac{z}{0} \stackrel{\text{def}}{=} \infty \quad (z \neq 0);
\]
другие операции с элементом $\infty$ не определены. $\varepsilon$-окрестностью точки $\infty$ называется множество
\[
U_{\varepsilon}(\infty) \stackrel{\text{def}}{=} \{z \in \mathbb{C} : |z| > 1/\varepsilon\} \cup \{\infty\} \equiv \hat{\mathbb{C}} \setminus \overline{U_{1/\varepsilon}(0)}.
\]

\textbf{6.1. Определение.} Множество $\hat{\mathbb{C}}$ с определённой выше топологией называется расширенной комплексной плоскостью.

\textbf{6.2. Теорема.} Топологические пространства $\hat{\mathbb{C}}$ и $S^2$ гомеоморфны.

\textbf{Доказательство.} (i) Рассмотрим сферу
\begin{equation}\tag{6.1}
	S^2 = \{(\xi, \eta, \zeta) \in \mathbb{R}^3 : \xi^2 + \eta^2 + \zeta^2 = 1\}
\end{equation}
и её северный полюс $N \stackrel{\text{def}}{=} (0, 0, 1)$.

Пусть $P = (\xi, \eta, \zeta)$ --- произвольная точка $S \setminus \{N\}$. Прямая $NP$ имеет параметрические уравнения
\[
\begin{cases}
	x = t\xi, \\
	y = t\eta, \\
	z = 1 + t(\zeta - 1).
\end{cases}
\]
Её пересечению $\pi(\xi, \eta, \zeta)$ с плоскостью $\mathbb{R}^2 \times \{0\}$ соответствует значение параметра $t = 1/(1 - \zeta)$:
\begin{equation}\tag{6.2}
	\pi(\xi, \eta, \zeta) = \left( \frac{\xi}{1 - \zeta}, \frac{\eta}{1 - \zeta}, 0 \right).
\end{equation}

Положим по определению
\[
\pi(N) \stackrel{\text{def}}{=} \infty,
\]
где точка $\infty$ определяется аналогично $\hat{\mathbb{C}}$; тогда ясно, что $\pi$ задаёт непрерывное отображение $S$ в $(\mathbb{R}^2 \times \{0\}) \cup \{\infty\}$.

(ii) Пусть $z = (x, y, 0) \in \mathbb{R}^2 \times \{0\}$ --- произвольная точка. Предположим, что $z = \pi(P)$ для некоторой точки $P = (\xi, \eta, \zeta) \in S \setminus \{N\}$.

Из формулы (6.2) мы имеем соотношения
\begin{equation}\tag{6.3}
	\xi = (1 - \zeta)x, \quad \eta = (1 - \zeta)y,
\end{equation}
подставляя которые в уравнение сферы (6.1), мы получим
\[
(1 - \zeta)^2 x^2 + (1 - \zeta)^2 y^2 + \zeta^2 = 1,
\]
или, после раскрытия скобок,
\[
(\|z\|^2 + 1)\zeta^2 - 2\|z\|^2\zeta + \|z\|^2 - 1 = 0.
\]
Это уравнение имеет корни
\[
\zeta = 1, \quad \zeta = \frac{\|z\|^2 - 1}{\|z\|^2 + 1}.
\]

первый из которых мы исключаем, поскольку $P \neq N$. Таким образом, отображение $\pi$ биективно; формулы для обратного отображения получаются подстановкой найденных значений $\zeta$ в (6.3):
\[
\pi^{-1}(z) = \left( \frac{2x}{\|z\|^2 + 1}, \frac{2y}{\|z\|^2 + 1}, \frac{\%z\|^2 - 1}{\|z\|^2 + 1} \right).
\]
Легко видеть, что это отображение непрерывно на $\mathbb{R}^2 \times \{0\}$, при этом $\lim_{z \to \infty} \pi^{-1}(z) = N$. Отождествляя $(\mathbb{R}^2 \times \{0\}) \cup \{\infty\}$ с $\hat{\mathbb{C}}$, мы получаем гомеоморфизм $S^2 \cong \hat{\mathbb{C}}$. $\quad \blacksquare$

\textbf{6.3. Определение.} Сферой Римана называется пара $(S^2, \pi)$, где $\pi$ --- гомеоморфизм сферы $S^2$ на $\hat{\mathbb{C}}$, называемый стереографической проекцией с центром в точке $(0,0,1)$ и определяемый формулой
\[
\pi(\xi, \eta, \zeta) = \frac{\xi + i\eta}{1 - \zeta}, \quad \pi(0,0,1) = \infty.
\]

\textbf{6.4. Определение.} Пусть $I \subset \mathbb{R}$ --- отрезок, и пусть $\varphi_1, \varphi_2 \in C^1(I, \hat{\mathbb{C}})$ --- пути в $\hat{\mathbb{C}}$, проходящие через $\infty$. Углом между $\varphi_1$ и $\varphi_2$ в точке $\infty$ называется угол между путями $1/\varphi_1$ и $1/\varphi_2$ в точке $0$ в предположении, что $(1/\varphi_1)'(0) \neq 0$, $(1/\varphi_2)'(0) \neq 0$.

\textbf{6.5. Определение.} Пусть $a \in \hat{\mathbb{C}}$. Функция $f : U(a) \to \hat{\mathbb{C}}$ называется голоморфной в точке $a$ соответственно при
\[
a = \infty, \ f(a) \neq \infty, \quad a \neq \infty, \ f(a) = \infty, \quad a = \infty, \ f(a) = \infty,
\]
если в точке $0$ голоморфна соответственно функция
\begin{equation}\tag{6.4}
	z \mapsto f(1/z), \quad z \mapsto 1/f(z), \quad z \mapsto 1/f(1/z).
\end{equation}
Мы говорим, что $f$ конформна в точке $a$, если соответствующая функция (6.4) конформна в нуле.

\textbf{6.6. Определение.} Биективное конформное отображение называется голоморфизмом. Голоморфизм множества на себя называется также автоморфизмом этого множества.

\textbf{6.7. Предложение.} Всякое дробно-линейное отображение
\begin{equation}\tag{6.5}
	w : z \mapsto \frac{az + b}{cz + d}, \quad ad \neq bc,
\end{equation}
где $w(\infty) \stackrel{\text{def}}{=} a/c$ и $w(-d/c) \stackrel{\text{def}}{=} \infty$, является автоморфизмом $\hat{\mathbb{C}}$.

\textbf{Доказательство.} В самом деле, всякое ДЛО (6.5) при $c \neq 0$ (случай $c = 0$ рассматривается тривиально) представимо в виде
\[
w(z) = \frac{a}{c} - \frac{ad - bc}{c(cz + d)} \stackrel{\text{def}}{=} A + \frac{B}{z + C}
\]
и потому является композицией
\begin{itemize}
	\item сдвига $z \mapsto z + C$,
	\item вращения с растяжением $z \mapsto Bz$ (заметим, что $B \neq 0$),
	\item инверсии относительно единичной окружности $z \mapsto 1/z$,
	\item сдвига $z \mapsto z + A$.
\end{itemize}
Легко видеть, что все перечисленные отображения являются автоморфизмами расширенной комплексной плоскости $\hat{\mathbb{C}}$. Поскольку композиция голоморфизмов есть голоморфизм, предложение доказано. $\quad \blacksquare$

%──────────────────────────────────────────────────────────────────────────────
\newpage
\section{Интеграл от функции комплексного переменного вдоль пути в $\mathbb{C}$. Его свойства}
Пусть даны путь $\phi \in C^1([\alpha, \beta], \mathbb{C})$ и функция $f \in C(\phi([\alpha, \beta]), \mathbb{C})$.

\textbf{Определение 7.1.} Интегралом от $f$ по пути $\phi$ называется число
\begin{equation}
	\int_{\phi} f \, dz \stackrel{\text{def}}{=} \int_{\alpha}^{\beta} f(\phi(t))\phi'(t) \, dt.
\end{equation}

\textbf{Пример 7.2.} Пусть $\phi(t) = a + r e^{it}$, $t \in [0, 2\pi]$, и $f(z) = (z - a)^n$. Тогда
\[
\int_{\phi} f \, dz = \int_{0}^{2\pi} r^n e^{int} \cdot ir e^{it} \, dt = ir^{n+1} \int_{0}^{2\pi} e^{i(n+1)t} \, dt =
\]
\[
= ir^{n+1} \int_{0}^{2\pi} \left[ \cos((n + 1)t) + i \sin((n + 1)t) \right] \, dt =
\begin{cases}
	0, & n \neq -1; \\
	2\pi i, & n = -1.
\end{cases} \quad \blacksquare
\]

\subsection{7.2 Свойства интеграла}

\textbf{Свойство 1.} Интеграл (7.1) является линейной комбинацией криволинейных интегралов 2-го рода.

\textit{Доказательство.} В самом деле, положим $f(x, y) = u(x, y) + iv(x, y)$, $\phi(t) = x(t) + iy(t)$. Тогда
\[
\int_{\phi} f \, dz = \int_{\alpha}^{\beta} f(\phi(t))\phi'(t) \, dt = \int_{\alpha}^{\beta} (ux' - vy') \, dt + i \int_{\alpha}^{\beta} (vx' + uy') \, dt =
\]
\[
= \int_{\phi} u \, dx - v \, dy + i \int_{\phi} v \, dx + u \, dy. \quad \blacksquare
\]

\textbf{Свойство 2.} Интеграл (7.1) линеен: если $f, g \in C(\phi([\alpha, \beta]), \mathbb{C})$ и $a, b \in \mathbb{C}$, то
\[
\int_{\phi} (af + bg) \, dz = a \int_{\phi} f \, dz + b \int_{\phi} g \, dz. \quad \blacksquare
\]

\textbf{Свойство 3. Ориентированность:} если $\phi^*(t) = \phi(\alpha + \beta - t)$, $t \in [\alpha, \beta]$, то
\[
\int_{\phi^*} f \, dz = - \int_{\phi} f \, dz. \quad \blacksquare
\]

\textbf{Свойство 4. Аддитивность:} если даны два пути $\phi_1 \in C^1([\alpha, \delta], \mathbb{C})$, $\phi_2 \in C^1([\delta, \beta], \mathbb{C})$ и $\phi = \phi_1 \cup \phi_2$, то 
\[
\int_{\phi} f \, dz = \int_{\phi_1} f \, dz + \int_{\phi_2} f \, dz. \quad \blacksquare
\]

\textbf{Свойство 5. Инвариантность относительно замены параметра:} если дан диффеоморфизм $\tau \in \text{Diff}^1([\alpha_1, \beta_1], [\alpha, \beta])$, то 
\[
\int_{\phi \circ \tau} f \, dz = \pm \int_{\phi} f \, dz,
\]
где знак выбирается в зависимости от типа монотонности $\tau$. Это обобщает свойство 3. $\quad \blacksquare$

\textbf{Свойство 6. Оценка модуля интеграла:}
\[
\left| \int_{\phi} f \, dz \right| \le \int_{\phi} |f| \, |dz|,
\]
где $|dz| = \sqrt{dx^2 + dy^2}$ — дифференциал длины кривой. В частности, если $\sup |f(z)| = M < +\infty$, то 
\[
\left| \int_{\phi} f \, dz \right| \le M L(\phi([\alpha, \beta])),
\]
где $L(\phi([\alpha, \beta]))$ — длина кривой $\phi([\alpha, \beta])$.

\textit{Доказательство.} В самом деле, пусть $\int_{\phi} f \, dz = I = |I|e^{i\theta}$; тогда
\[
|I| = I e^{-i\theta} = \int_{\alpha}^{\beta} e^{-i\theta} f(\phi(t))\phi'(t) \, dt \stackrel{|I| \in \mathbb{R}}{\equiv} \int_{\alpha}^{\beta} \text{Re} \left[ e^{-i\theta} f(\phi(t))\phi'(t) \right] \, dt \le
\]
\[
\le \int_{\alpha}^{\beta} \left| e^{-i\theta} f(\phi(t))\phi'(t) \right| \, dt = \int_{\alpha}^{\beta} |f(\phi(t))| |\phi'(t)| \, dt \equiv \int_{\phi} |f| \, |dz|. \quad \blacksquare
\]

%──────────────────────────────────────────────────────────────────────────────
\newpage
\section{Теорема Коши для односвязных и многосвязных областей}
\textbf{Теорема 3 (Коши для односвязной области).} Если функция $f$ голоморфна в односвязной области $D \subset \mathbb{C}$, то ее интеграл вдоль любого замкнутого пути $\gamma \subset D$ равен нулю.

Ввиду важности этой теоремы приведем еще ее элементарное доказательство в двух дополнительных предположениях: 1) производная $f'$ непрерывна в $D$ и 2) $\gamma$ — гладкий жорданов путь.

Из второго предположения следует, что $\gamma$ является границей области $G$, принадлежащей области $D$ в силу односвязности последней. Первое предположение позволяет применить известную из анализа формулу Римана --- Грина
\begin{equation}
	\int_{\partial G} P \, dx + Q \, dy = \iint_{G} \left( \frac{\partial Q}{\partial x} - \frac{\partial P}{\partial y} \right) dx \, dy,
\end{equation}
при выводе которой требуется непрерывность в $\overline{G}$ частных производных функций $P$ и $Q$ (через $\partial G$ здесь обозначается граница области $G$, проходимая против часовой стрелки). Применяя эту формулу к действительной и мнимой частям интеграла
\[
\int_{\partial G} f \, dz = \int_{\partial G} u \, dx - v \, dy + i \int_{\partial G} v \, dx + u \, dy,
\]
мы получим
\[
\int_{\partial G} f \, dz = \iint_{G} \left\{ -\frac{\partial v}{\partial x} - \frac{\partial u}{\partial y} + i \left( \frac{\partial u}{\partial x} - \frac{\partial v}{\partial y} \right) \right\} dx \, dy.
\]
Пользуясь символом комплексной производной $\frac{\partial}{\partial \bar{z}}$ (см. п. 6), мы перепишем последнее соотношение в виде формулы
\begin{equation}
	\int_{\partial G} f \, dz = 2i \iint_{G} \frac{\partial f}{\partial \bar{z}} \, dx \, dy,
\end{equation}
которую можно рассматривать как комплексную запись формулы Римана --- Грина.

Так как в силу голоморфности $\frac{\partial f}{\partial \bar{z}} \equiv 0$, то теорема Коши (в сделанных дополнительных предположениях) непосредственно вытекает из этой формулы. $\quad \blacksquare$

\textbf{Теорема 1. (Коши для многосвязных областей)} Пусть функция $f \in H(D)$ и $G$ — любая компактно принадлежащая $D$ область, ограниченная конечным числом кривых. Тогда интеграл от $f$ по ориентированной границе области $G$ равен нулю:
\begin{equation}
	\int_{\partial G} f \, dz = 0.
\end{equation}

\textbf{Доказательство.} Проведем в $G$ конечное число разрезов $\lambda_\nu^{\pm}$, связывающих компоненты границы этой области (на рис. 32 мы для наглядности изобразим эти разрезы, состоящими из двух берегов). Очевидно, что замкнутая кривая $\Gamma$, состоящая из ориентированной границы $\partial G$ и совокупностей $\Lambda^+ = \bigcup \lambda_\nu^+$ и $\Lambda^- = \bigcup \lambda_\nu^-$, гомотопна нулю в области $D$. В силу свойств интегралов
\[
\int_{\Gamma} f \, dz = \int_{\partial G} f \, dz + \int_{\Lambda^+} f \, dz + \int_{\Lambda^-} f \, dz = \int_{\partial G} f \, dz,
\]
а так как $\Gamma \sim 0$ в $D$, то по интегральной теореме Коши для контура, гомотопного нулю, интеграл по $\Gamma$ равен нулю. Отсюда мы получаем (1):
\[
\int_{\partial G} f \, dz = 0. \quad \blacksquare
\]

%──────────────────────────────────────────────────────────────────────────────
\newpage
\section{Интегральная формула Коши для функции и ее производных}
\textbf{Теорема 8.5 (Интегральная формула Коши).} Пусть $D$ — односвязная область в $\mathbb{C}$ и $f \in C^{\omega}(D, \mathbb{C})$. Тогда для любой $z_0 \in D$ выполнено равенство
\begin{equation}
	f(z_0) = \frac{1}{2\pi i} \oint_{\partial D} \frac{f(z)}{z - z_0} \, dz.
\end{equation}

\textbf{Доказательство.} Рассмотрим окрестность $U_{\delta}(z_0) \subset D$ и окружность $L_1 = \partial U_{\delta}(z_0)$. Функция 
\[
z \mapsto \frac{f(z)}{z - z_0}
\]
голоморфна в области $D \setminus U_{\delta}(z_0)$. По теореме Коши для многосвязной области имеем:
\[
0 = \oint_{\partial D - L_1} \frac{f(z)}{z - z_0} \, dz = \oint_{\partial D} \frac{f(z)}{z - z_0} \, dz - \oint_{L_1} \frac{f(z)}{z - z_0} \, dz.
\]
Далее, используя пример 7.2 и свойство об оценке модуля, получаем:
\[
\left| \oint_{L_1} \frac{f(z)}{z - z_0} \, dz - 2\pi i f(z_0) \right| = \left| \oint_{L_1} \frac{f(z)}{z - z_0} \, dz - \oint_{L_1} \frac{f(z_0)}{z - z_0} \, dz \right| =
\]
\[
= \left| \oint_{L_1} \frac{f(z) - f(z_0)}{z - z_0} \, dz \right| \le \oint_{L_1} \left| \frac{f(z) - f(z_0)}{z - z_0} \right| \, |dz| < \frac{\varepsilon}{\delta} \cdot 2\pi\delta = 2\pi\varepsilon,
\]
если $|f(z) - f(z_0)| < \varepsilon$. В силу произвольности $\varepsilon$ имеем:
\[
\oint_{L_1} \frac{f(z)}{z - z_0} \, dz = 2\pi i f(z_0). \quad \blacksquare
\]

%──────────────────────────────────────────────────────────────────────────────
\newpage
\section{Степенные ряды в $\mathbb{C}$, их свойства. Голоморфность суммы степенного ряда}

Пусть $M \subset \hat{\mathbb{C}}$ и $f_n : M \to \mathbb{C}$ для $n \in \mathbb{N}$. Рассмотрим ряд
\begin{equation}
	\sum_{n=1}^{\infty} f_n.
\end{equation}
Повторим некоторые свойства функциональных рядов. Их доказательство аналогично действительному случаю.

\textbf{Теорема 9.1 (Признак Вейерштрасса).} Если ряд (1) сходится нормально, то он сходится равномерно.

\textit{Доказательство} аналогично вещественному случаю и сразу вытекает из критерия Коши (для числового и функционального рядов) и неравенства треугольника. $\quad \blacksquare$

Повторим некоторые свойства степенных рядов вида
\begin{equation}
	\sum_{n=0}^{\infty} a_n(z - z_0)^n.
\end{equation}
Их доказательство аналогично действительному случаю.

\textbf{Теорема 9.2 (Первая теорема Абеля).} Если ряд (2) сходится в точке $z_1 \in \mathbb{C}$, то он сходится абсолютно и равномерно на любом множестве $K \subset\subset \{z \in \mathbb{C} : |z - z_0| < |z_1 - z_0|\}$. $\quad \blacksquare$

\textbf{Теорема 9.3 (Формула Коши — Адамара).} Пусть
\begin{equation}
	R = \frac{1}{\lim\limits_{n \to \infty} \sqrt[n]{|a_n|}}.
\end{equation}
Тогда ряд (2) сходится абсолютно при $|z - z_0| < R$ и расходится при $|z - z_0| > R$. $\quad \blacksquare$

В частности, если ряд (2) сходится условно в точке $z_1 \in \mathbb{C}$, то $R = |z_1 - z_0|$. Число (3) называется \textit{радиусом сходимости} ряда (2).

Теорема 5. \textit{Сумма степенного ряда}
\begin{equation}
	f(z) = \sum_{n=0}^{\infty} c_n (z - a)^n
	\tag{12}
\end{equation}
\textit{голоморфна в круге его сходимости.}

$\blacktriangleleft$ Мы предполагаем, что радиус сходимости ряда $R > 0$, ибо иначе нечего доказывать. Напишем формально производный ряд
\begin{equation}
	\sum_{n=1}^{\infty} n c_n (z - a)^{n-1} = \varphi(z);
	\tag{13}
\end{equation}
он сходится или расходится вместе с рядом $\sum_{n=1}^{\infty} n c_n (z - a)^n$, а так как
$$
\varlimsup_{n \to \infty} \sqrt[n]{n |c_n|} = \varlimsup_{n \to \infty} \sqrt[n]{|c_n|},
$$
то радиус сходимости ряда (13) тоже равен $R$. На компактных подмножествах круг $U = \{|z - a| < R\}$ ряд (13) сходится равномерно, следовательно, функция $\varphi(z)$ непрерывна в этом круге.

По той же причине ряд (13) можно почленно интегрировать по границе любого треугольника $\Delta \Subset U$:
$$
\int_{\partial \Delta} \varphi \, dz = \sum_{n=1}^{\infty} n c_n \int_{\partial \Delta} (z - a)^{n-1} \, dz = 0
$$
(по теореме Коши все интегралы в правой части равны нулю, значит, и интеграл в левой части также). Следовательно, можно применить лемму 1 из п. 15, по которой функция
$$
\int_{[a, z]} \varphi(\zeta) \, d\zeta = \sum_{n=1}^{\infty} n c_n \int_{[a, z]} (\zeta - a)^{n-1} \, d\zeta = \sum_{n=1}^{\infty} c_n (z - a)^n
$$

(мы снова воспользовались равномерной сходимостью) в каждой точке $z \in U$ имеет производную, равную $\varphi(z)$. Но тогда и функция
$$
f(z) = c_0 + \int_{[a, z]} \varphi(\zeta) \, d\zeta
$$
в каждой точке $z \in U$ имеет производную $f'(z) = \varphi(z) \blacktriangleright$

%──────────────────────────────────────────────────────────────────────────────
\newpage
\section{Теорема о разложении голоморфной функции в ряд Тейлора. Неравенства Коши для коэффициентов ряда Тейлора. Теорема Лиувилля}

\subsection{9.2 Теорема о разложении голоморфной функции в степенной ряд}

\textbf{Теорема 9.4.} Пусть $D \subset \mathbb{C}$ — открытое множество и $f \in C^{\omega}(D, \mathbb{C})$. Тогда в любом круге $U = \{z \in \mathbb{C} : |z - z_0| < R\} \subset D$, где $z_0 \in D$ и $R > 0$, функция $f$ допускает представление
\begin{equation}
	f(z) = \sum_{n=0}^{\infty} a_n(z - z_0)^n,
\end{equation}
где коэффициенты $a_n$ вычисляются по формуле
\[
a_n = \frac{1}{2\pi i} \oint_{\partial U} \frac{f(z)}{(z - z_0)^{n+1}} \, dz, \quad n \in \mathbb{N} \cup \{0\}.
\]

\textbf{Доказательство.} Выберем любую точку $z \in U$ и число $r$ так, что $|z - z_0| < r < R$. По формуле Коши:
\begin{equation}
	f(z) = \frac{1}{2\pi i} \oint_{\partial U} \frac{f(\xi)}{\xi - z} \, d\xi.
\end{equation}
При этом, если $\xi \in \partial U$, то $\left| \frac{z - z_0}{\xi - z_0} \right| < 1$ и
\[
\frac{1}{\xi - z} = \frac{1}{\xi - z_0 - (z - z_0)} = \frac{1}{\xi - z_0} \cdot \frac{1}{1 - \dfrac{z - z_0}{\xi - z_0}} = \sum_{n=0}^{\infty} \frac{(z - z_0)^n}{(\xi - z_0)^{n+1}}.
\]
Подставим это выражение в интеграл (5) и произведем почленное интегрирование:
\[
f(z) = \frac{1}{2\pi i} \oint_{\partial U} f(\xi) \sum_{n=0}^{\infty} \left[ \frac{(z - z_0)^n}{(\xi - z_0)^{n+1}} \right] \, d\xi = \frac{1}{2\pi i} \sum_{n=0}^{\infty} \left[ \oint_{\partial U} \frac{f(\xi)}{(\xi - z_0)^{n+1}} \, d\xi \right] (z - z_0)^n.
\]
Корректность этого перехода следует из того, что функция $f$ непрерывна на $D$, а потому ограничена на $\partial U$ (по признаку Вейерштрасса полученный функциональный ряд сходится равномерно на контуре интегрирования). Требуемое на этом доказано. $\quad \blacksquare$

Ряд (4) называется \textit{рядом Тейлора} функции $f$.

\subsection{Неравенства Коши}

\textbf{Предложение 9.5 (Неравенства Коши для коэффициентов ряда Тейлора).} Пусть даны замкнутый круг $[U] = \{z \in \mathbb{C} : |z - z_0| \le r\} \subset D$ и функция $f \in C^{\omega}([U], \mathbb{C})$. Обозначим $M = \sup\limits_{z \in [U]} |f(z)|$. Тогда для коэффициентов $a_n$ ряда Тейлора в $[U]$ справедливы неравенства:
\begin{equation}
	|a_n| \le \frac{M}{r^n}, \quad n \in \mathbb{N} \cup \{0\}.
\end{equation}

\textbf{Доказательство.} Из Теоремы о разложении голоморфной функции в ряд Тейлора известно интегральное выражение для коэффициентов ряда Тейлора по границе круга $\partial U = \{z \in \mathbb{C} : |z - z_0| = r\}$:
\[
a_n = \frac{1}{2\pi i} \oint_{\partial U} \frac{f(z)}{(z - z_0)^{n+1}} \, dz.
\]
Применим свойство оценки модуля интеграла. Для любой точки $z \in \partial U$ имеем $|z - z_0| = r$ и $|f(z)| \le M$. Длина окружности интегрирования равна $L(\partial U) = 2\pi r$. Получаем:
\[
|a_n| = \left| \frac{1}{2\pi i} \oint_{\partial U} \frac{f(z)}{(z - z_0)^{n+1}} \, dz \right| \le \frac{1}{2\pi} \oint_{\partial U} \frac{|f(z)|}{|z - z_0)^{n+1}|} \, |dz| \le
\]
\[
\le \frac{1}{2\pi} \cdot \frac{M}{r^{n+1}} \oint_{\partial U} |dz| = \frac{1}{2\pi} \cdot \frac{M}{r^{n+1}} \cdot 2\pi r = \frac{M}{r^n}. \quad \blacksquare
\]

\subsection{Теорема Лиувилля}

\textbf{Теорема Лиувилля.} Если функция $f(z)$ голоморфна во всей комплексной плоскости $\mathbb{C}$ (является целой функцией) и ограничена по модулю, то она является константой.

\textbf{Доказательство.} 
По условию функция $f(z)$ ограничена на $\mathbb{C}$. Это означает, что существует такая вещественная постоянная $M > 0$, что для всех $z \in \mathbb{C}$ выполняется неравенство:
\begin{equation}
	|f(z)| \le M.
\end{equation}

Поскольку функция $f(z)$ голоморфна во всей плоскости $\mathbb{C}$, её можно разложить в ряд Тейлора с центром в любой точке (для удобства выберем начало координат $z_0 = 0$), причем радиус сходимости этого ряда равен бесконечности ($R = \infty$):
\[
f(z) = \sum_{n=0}^{\infty} a_n z^n.
\]

Рассмотрим произвольный замкнутый круг $[U_r] = \{z \in \mathbb{C} : |z| \le r\}$ некоторого радиуса $r > 0$. Согласно \textbf{неравенствам Коши для коэффициентов ряда Тейлора}, для любого индекса $n \in \mathbb{N}$ справедливо соотношение:
\begin{equation}
	|a_n| \le \frac{\sup_{|z|=r} |f(z)|}{r^n}.
\end{equation}

Используя условие ограниченности (1), мы можем усилить неравенство (2) для любого фиксированного $r$:
\begin{equation}
	|a_n| \le \frac{M}{r^n}.
\end{equation}

Поскольку функция $f(z)$ голоморфна везде, мы имеем право устремить радиус круга $r$ к бесконечности. Рассмотрим предел правой части неравенства (3) при $r \to \infty$ для всех $n \ge 1$:
\[
\lim_{r \to \infty} \frac{M}{r^n} = 0 \quad (\text{так как } n \ge 1).
\]
Следовательно, для любого $n \in \{1, 2, 3, \dots\}$ выполняется:
\[
|a_n| = 0 \implies a_n = 0.
\]

Таким образом, все коэффициенты ряда Тейлора, начиная с первого, строго равны нулю. Подставляя их в разложение функции, получаем:
\[
f(z) = a_0 + 0 \cdot z + 0 \cdot z^2 + \dots \equiv a_0.
\]
Так как $a_0 = f(0) = \text{const}$, то $f(z) \equiv \text{const}$ на всей комплексной плоскости $\mathbb{C}$. Теорема доказана. $\quad \blacksquare$

%──────────────────────────────────────────────────────────────────────────────
\newpage
\section{Бесконечная дифференцируемость голоморфных функций. Единственность разложения в степенной ряд. Теорема Мореры. Эквивалентность голоморфности в смысле Римана, Коши и Вейерштрасса}

Теорема 1. \textit{Производная функции $f \in H(D)$ голоморфна в области $D$.}

\textbf{Доказательство:}
$\blacktriangleleft$ Для любой точки $z_0 \in D$ мы построим круг $U = \{|z - z_0| < R\}$, принадлежащий $D$. По теореме 1 п. 19 функция $f$ в этом круге представляется как сумма степенного ряда. По теореме 5 п. 19 производная $f' = \varphi$ представляется рядом, сходящимся в том же круге. Поэтому к $\varphi$ можно снова применить теорему 5, и, значит, $\varphi$ дифференцируема в $U$ в смысле комплексного анализа $\blacktriangleright$

Теорема 2. \textit{Если функция $f$ в круге $\{|z - z_0| < R\}$ представима как сумма степенного ряда}
\begin{equation}
	f(z) = \sum_{n=0}^{\infty} c_n (z - z_0)^n,
	\tag{1}
\end{equation}
\textit{то коэффициенты этого ряда определяются однозначно по формулам}
\begin{equation}
	c_n = \frac{f^{(n)}(z_0)}{n!} \quad (n = 0, 1, \dots).
	\tag{2}
\end{equation}

$\blacktriangleleft$ Подставляя в (1) $z = z_0$, найдем $f(z_0) = c_0$. Дифференцируя ряд (1) почленно:
$$
f'(z) = c_1 + 2c_2(z - z_0) + 3c_3(z - z_0)^2 + \dots,
$$
и затем подставляя $z = z_0$, найдем $f'(z_0) = c_1$. Продифференцируем (1) $n$ раз:
$$
f^{(n)}(z) = n! \, c_n + c'_1 (z - z_0) + c'_2 (z - z_0)^2 + \dots
$$
(мы не выписываем выражений для коэффициентов), и снова подставим $z = z_0$; получим $n! \, c_n = f^{(n)}(z_0) \blacktriangleright$

\textbf{Теорема 3 (Морера).} Если функция $f$ непрерывна в области $D$ и интеграл от нее по границе $\partial \Delta$ любого треугольника $\Delta \Subset D$ равен нулю, то $f \in H(D)$.

\textbf{Доказательство:}
$\blacktriangleleft$ Для любой точки $a \in D$ построим круг $U = \{|z - a| < r\} \subset D$. По теореме 2 п. 15 функция
$$
F(z) = \int_{[a, z]} f(\zeta) \, d\zeta
$$
голоморфна в $U$ и $F'(z) = f(z)$ в каждой точке $z \in U$. По теореме 1 отсюда следует, что и $f$ голоморфна в $U$. Тем самым доказана голоморфность $f$ в каждой точке $a \in D \blacktriangleright$

Теорема 4. \textit{Эквивалентные следующие три утверждения:}

(R) \textit{функция $f$ в некоторой окрестности $U$ точки $a$ имеет производную $f'(z)$ в смысле комплексного анализа;}

(C) \textit{функция $f$ непрерывна в некоторой окрестности $U$ точки $a$, и интеграл от нее по границе любого треугольника $\Delta \Subset U$ равен нулю;}

(W) \textit{функция $f$ разлагается в степенной ряд, сходящийся в некоторой окрестности $U$ точки $a$.}

Эти три утверждения отражают три концепции в построении теории голоморфных функций. Обычно функцию, удовлетворяющую условию (R), называют голоморфной в смысле Римана, условию (C) — голоморфной в смысле Коши и условию (W) — голоморфной в смысле Вейерштрасса.

Импликация $(R) \Rightarrow (C)$ доказана в теореме Коши (п. 16), $(C) \Rightarrow (R)$ — в теореме Мореры, $(R) \Rightarrow (W)$ — в теореме Тейлора, $(W) \Rightarrow (R)$ — в теореме о голоморфности суммы степенного ряда.

%──────────────────────────────────────────────────────────────────────────────
\newpage
\section{Нули голоморфной функции, их свойства. Теорема единственности. Вычисление порядка нуля}
\textit{Нулем функции $f$ называется любая точка $a \in \mathbb{C}$, в которой эта функция равна нулю, т. е. корень уравнения $f(z) = 0$.}

Теорема 1. \textit{Если точка $a$ является нулем голоморфной в этой точке функции $f$, не равной тождественно нулю ни в} какой окрестности $a$, то существует такое натуральное число $n$, что
\begin{equation}
	f(z) = (z - a)^n \varphi(z),
	\tag{1}
\end{equation}
где функция $\varphi$ голоморфна в точке $a$ и отлична от нуля в некоторой окрестности этой точки.

$\blacktriangleleft$ В самом деле, в некоторой окрестности точки $a$ функция $f$ разлагается в степенной ряд. Свободный член этого ряда равен нулю, ибо $f(a) = 0$, но все его коэффициенты не могут быть равными нулю, ибо тогда $f \equiv 0$ в некоторой окрестности $a$. Поэтому найдется отличный от нуля коэффициент с младшим номером, который мы обозначим через $n$, и разложение имеет вид
\begin{equation}
	f(z) = c_n (z - a)^n + c_{n+1} (z - a)^{n+1} + \dots, \quad c_n \neq 0.
	\tag{2}
\end{equation}
Обозначим через
$$
\varphi(z) = c_n + c_{n+1} (z - a) + \dots
$$
сумму ряда, сходящегося в некоторой окрестности $a$, и поэтому голоморфную в этой окрестности. Так как $\varphi(a) = c_n \neq 0$, то в силу непрерывности этой функции $\varphi \neq 0$ в некоторой окрестности $a \blacktriangleright$

Теорема 2 (единственности). \textit{Если две функции $f_1, f_2 \in H(D)$ совпадают на множестве $E$, которое имеет хотя бы одну предельную точку $a$, принадлежащую $D$, то $f_1 = f_2$ всюду в $D$.}

$\blacktriangleleft$ Функция $f = f_1 - f_2 \in H(D)$; нужно доказать, что $f \equiv 0$ в $D$, т. е. что множество $\mathcal{F} = \{z \in D : f(z) = 0\} \supset E$ совпадает с $D$. Предельная точка $a$ является нулем $f$ (в силу непрерывности последней). По теореме 1 функция $f = 0$ в некоторой окрестности $a$, ибо иначе эта точка не могла быть предельной для множества нулей $f$.

Таким образом, открытое ядро $\mathring{\mathcal{F}}$ множества $\mathcal{F}$ (т. е. совокупность его внутренних точек) непусто — оно содержит точку $a$. По построению $\mathring{\mathcal{F}}$ открыто, но оно в то же время и замкнуто (в относительной топологии области $D$). В самом деле, если $b \in D$ является предельной точкой $\mathring{\mathcal{F}}$, то по той же теореме 1 функция $f \equiv 0$ в некоторой окрестности $b$, т. е. $b \in \mathring{\mathcal{F}}$. Так как $D$ по определению области связно, то по теореме 3 п. 4 имеем $\mathring{\mathcal{F}} = D \blacktriangleright$

\textbf{Определение 10.5.} Порядком нуля $a \in \mathbb{C}$ голоморфной в точке $a$ функции $f$ называется число
$$
n = \min\{k \in \mathbb{N} : f^{(k)}(a) \neq 0\}.
$$
\textbf{Теорема 10.6.} Порядок нуля $a \in \mathbb{C}$ голоморфной в точке $a$ функции $f$ совпадает с максимальной степенью бинома $(z - a)^n$, при которой частное $\frac{f(z)}{(z - a)^n}$ голоморфно в $a$.

\textbf{Доказательство.} Это сразу следует из теоремы 1. $\blacksquare$

%──────────────────────────────────────────────────────────────────────────────
\newpage
\section{Ряды Лорана, их области сходимости. Теоремы о разложении голоморфной функции и о единственности разложения в ряд Лорана. Неравенства Коши для коэффициентов ряда Лорана}
Рассмотрим ряд
\begin{equation}
	\sum_{n=-\infty}^{\infty} c_n (z - a)^n \stackrel{\text{def}}{=} \sum_{n=0}^{\infty} c_n (z - a)^n + \sum_{n=1}^{\infty} c_{-n} (z - a)^{-n}.
	\tag{10.1}
\end{equation}
Легко видеть, что ряды в правой части сходятся соответственно при
$$
|z - a| < R \equiv \frac{1}{\varlimsup_{n \to \infty} \sqrt[n]{|c_n|}},
$$
$$
|z - a| > r \equiv \varlimsup_{n \to \infty} \sqrt[n]{|c_{-n}|}.
$$
Следовательно, областью сходимости ряда (10.1) в общем случае является кольцо
$$
r < |z - a| < R.
$$

\textbf{Теорема 1 (Лоран).} \textit{Любую функцию $f$, голоморфную в кольце $V = \{r < |z - a| < R\}$, можно в этом кольце представить как сумму сходящегося ряда}
\begin{equation}
	f(z) = \sum_{n=-\infty}^{\infty} c_n (z - a)^n,
	\tag{1}
\end{equation}
\textit{коэффициенты которого определяются по формулам}
\begin{equation}
	c_n = \frac{1}{2\pi i} \int_{\{|z - a| = \rho\}} \frac{f(\zeta) \, d\zeta}{(\zeta - a)^{n+1}} \quad (n = 0, \pm 1, \pm 2, \dots),
	\tag{2}
\end{equation}
\textit{где $r < \rho < R$.}

$\blacktriangleleft$ Фиксируем произвольную точку $z \in V$ и построим кольцо $V' = \{\zeta : r' < |\zeta - a| < R'\}$ такое, что $z \in V' \Subset V$. По интегральной формуле Коши имеем
\begin{equation}
	f(z) = \frac{1}{2\pi i} \int_{\partial V'} \frac{f(\zeta) \, d\zeta}{\zeta - z} = \frac{1}{2\pi i} \int_{\Gamma'} \frac{f(\zeta) \, d\zeta}{\zeta - z} - \frac{1}{2\pi i} \int_{\gamma'} \frac{f(\zeta) \, d\zeta}{\zeta - z},
	\tag{3}
\end{equation}
где окружности $\Gamma' = \{|\zeta - a| = R'\}$ и $\gamma' = \{|\zeta - a| = r'\}$ ориентированы против часовой стрелки.

Для всех $\zeta \in \Gamma'$ имеем $\left| \frac{z - a}{\zeta - a} \right| = q < 1$, поэтому геометрическая прогрессия
$$
\frac{1}{\zeta - z} = \frac{1}{(\zeta - a)\left(1 - \frac{z - a}{\zeta - a}\right)} = \sum_{n=0}^{\infty} \frac{(z - a)^n}{(\zeta - a)^{n+1}}
$$
сходится равномерно по $\zeta$ на $\Gamma'$. Умножая ее на ограниченную функцию $\frac{1}{2\pi i} f(\zeta)$ (что не нарушает равномерной сходимости) и интегрируя почленно вдоль $\Gamma'$, получаем
\begin{equation}
	\frac{1}{2\pi i} \int_{\Gamma'} \frac{f(\zeta) \, d\zeta}{\zeta - z} = \sum_{n=0}^{\infty} c_n (z - a)^n,
	\tag{4}
\end{equation}
где
\begin{equation}
	c_n = \frac{1}{2\pi i} \int_{\Gamma'} \frac{f(\zeta) \, d\zeta}{(\zeta - a)^{n+1}} \quad (n = 0, 1, \dots).
	\tag{5}
\end{equation}

Второй интеграл в формуле (3) придется разлагать иначе. При всех $\zeta \in \gamma'$ имеем $\left| \frac{\zeta - a}{z - a} \right| = q_1 < 1$, поэтому мы получаем равномерно сходящуюся на $\gamma'$ геометрическую прогрессию так:
$$
-\frac{1}{\zeta - z} = \frac{1}{(z - a)\left(1 - \frac{\zeta - a}{z - a}\right)} = \sum_{n=1}^{\infty} \frac{(\zeta - a)^{n-1}}{(z - a)^n}.
$$
Снова умножая ее на ограниченную функцию $\frac{1}{2\pi i} f(\zeta)$ и интегрируя почленно вдоль $\gamma'$, получаем
\begin{equation}
	-\frac{1}{2\pi i} \int_{\gamma'} \frac{f(\zeta) \, d\zeta}{\zeta - z} = \sum_{n=1}^{\infty} \frac{d_n}{(z - a)^n},
	\tag{6}
\end{equation}
где
\begin{equation}
	d_n = \frac{1}{2\pi i} \int_{\gamma'} f(\zeta) (\zeta - a)^{n-1} \, d\zeta \quad (n = 1, 2, \dots).
	\tag{7}
\end{equation}

Заменим теперь в формулах (6) и (7) индекс $n$, пробегающий значения $1, 2, \dots$, индексом $-n$, пробегающим значения $-1, -2, \dots$ (это ничего не меняет))
\begin{equation}
	c_n = d_{-n} = \frac{1}{2\pi i} \int_{\gamma'} f(\zeta) (\zeta - a)^{-n-1} \, d\zeta \quad (n = -1, -2, \dots);
	\tag{8}
\end{equation}
тогда разложение (6) примет вид
\begin{equation}
	-\frac{1}{2\pi i} \int_{\gamma'} \frac{f(\zeta) \, d\zeta}{\zeta - z} = \sum_{n=-\infty}^{-1} c_n (z - a)^n.
	\tag{6'}
\end{equation}

Теперь подставим (4) и (6') в (3); получим нужное разложение (1):
$$
f(z) = \sum_{n=-\infty}^{\infty} c_n (z - a)^n,
$$
где ряд определяется как объединение рядов (4) и (6'). Остается заметить, что по теореме об инвариантности интеграла в формулах (5) и (8) окружности $\Gamma'$ и $\gamma'$ можно заменить любой окружностью $\{|\zeta - a| = \rho\}$, где $r < \rho < R$, и тогда эти формулы примут вид (2) $\blacktriangleright$

\textbf{Неравенства Коши} (для коэффициентов ряда Лорана). 
Пусть функция $f$ голоморфна в кольце $V = \{ r < |z - a| < R \}$ и 
на окружности $\gamma_\rho = \{ |z - a| = \rho \}$, $r < \rho < R$, ее модуль не превосходит 
постоянной $M$. Тогда коэффициенты ряда Лорана функции $f$ в кольце $V$ 
удовлетворяют неравенствам
\begin{equation}
	|c_n| \leqslant \frac{M}{\rho^n} \qquad (n = 0, \pm 1, \pm 2, \ldots).
\end{equation}

%──────────────────────────────────────────────────────────────────────────────
\newpage
\section{Изолированные особые точки голоморфных функций, их классификация и характеризация в терминах рядов Лорана. Поведение голоморфных функций в окрестности особых точек}

\textbf{Определение 1.} Точка $a \in \mathbb{C}$ называется \textit{изолированной особой точкой} функции $f$, если существует такая проколотая окрестность этой точки (т. е. множество $\{0 < |z - a| < r\}$, если точка $a$ конечна, или множество $\{R < |z| < \infty\}$, если $a = \infty$), в которой функция $f$ голоморфна.

В зависимости от поведения $f$ при приближении к такой точке различают три типа особых точек.

\textbf{Определение 2.} Изолированная особая точка $a$ функции $f$ называется
\begin{enumerate}
	\item[(\text{I})] \textit{устранимой точкой}, если существует конечный
	$$\lim_{z \to a} f(z) = A;$$
	\item[(\text{II})] \textit{полюсом}, если существует
	$$\lim_{z \to a} f(z) = \infty;$$
	\item[(\text{III})] \textit{существенно особой точкой}, если $f$ не имеет ни конечного, ни бесконечного предела при $z \to a$.
\end{enumerate}

\textbf{Теорема 1} (характеризация устранимых особых точек). Следующие условия эквивалентны:
\begin{enumerate}
	\item[(\text{I})] $a$ --- устранимая особая точка $f$;
	\item[(\text{II})] последовательность $(c_{-n})$ нулевая;
	\item[(\text{III})] $f$ ограничена в некоторой окрестности точки $a$.
\end{enumerate}

\textbf{Доказательство.} $(\text{I}) \implies (\text{III})$ Это стандартное свойство предела функции.

$(\text{III}) \implies (\text{II})$ Пусть $\mathring{V}(a) \subset \mathring{U}(a)$ и $|f(z)| \le M$ в $\mathring{V}(a)$. Рассмотрим некоторую окружность $\{|z - a| = \rho\} \subset \mathring{V}(a)$. Тогда справедливы неравенства Коши для коэффициентов ряда Лорана
$$|c_n| \le \frac{M}{\rho^n}, n \in -\mathbb{N}^\times,$$
из которых в пределе при $\rho \to 0$ получаем, что $c_n = 0$ при $n < 0$.

$(\text{II}) \implies (\text{I})$ Если
$$f(z) = \sum_{n=0}^{\infty} c_n (z - a)^n, \quad z \in \mathring{U}(a),$$
то почленный переход к пределу даёт
$$\lim_{z \to a} f(z) = c_0,$$
так что $a$ --- устранимая особая точка. $\blacksquare$

\textbf{Теорема 2} (характеризация полюса). Следующие условия эквивалентны:
\begin{enumerate}
	\item[(\text{I})] $a$ --- полюс $f$;
	\item[(\text{II})] последовательность $(c_{-n})$ ненулевая и финитная;
	\item[(\text{III})] точка $a$ является изолированной особой точкой функции $\varphi = 1/f$ и $\lim_{z \to a} \varphi(z) = 0$.
\end{enumerate}

\textbf{Доказательство.} $(\text{I}) \implies (\text{III})$ Это очевидно.

$(\text{III}) \implies (\text{II})$ Доопределим $\varphi$ нулём в точке $a$. Тогда $\varphi$ голоморфна в окрестности $a$, и её разложение в ряд Тейлора по степеням $z - a$ имеет вид
$$\varphi(z) = \sum_{n=N}^{\infty} b_n (z - a)^n.$$
Тогда
$$f(z) = \frac{1}{\varphi(z)} = \frac{1}{(z - a)^N} \cdot \frac{1}{\sum_{n=N}^{\infty} b_n (z - a)^{n-N}} \overset{\text{def}}{=} \frac{1}{(z - a)^N} \sum_{n=0}^{\infty} a_n (z - a)^n$$
(здесь мы разложили вторую дробь в ряд Тейлора, воспользовавшись её голоморфностью). Это означает выполнение $(\text{II})$.

$(\text{II}) \implies (\text{I})$ Если
$$f(z) = \sum_{n=-N}^{\infty} c_n (z - a)^n, \quad z \in \mathring{U}(a),$$
то почленный переход к пределу даёт
$$\lim_{z \to a} f(z) = \infty,$$
так что $a$ --- полюс. $\blacksquare$

\textbf{Теорема 3.} Изолированная особая точка $a$ функции $f$ является существенно особой в том и только том случае, если главная часть лорановского разложения $f$ в окрестности точки $a$ содержит бесконечно много отличных от нуля членов.

$\blacktriangleleft$ Теорема, по существу, уже содержится в теоремах 1 и 2 (если главная часть содержит бесконечное число членов, то $a$ не может быть ни устранимой точкой, ни полюсом; если $a$ --- существенно особая точка, то главная часть не может ни отсутствовать, ни содержать конечное число членов) $\blacktriangleright$

%──────────────────────────────────────────────────────────────────────────────
\newpage
\section{Вычеты, их вычисление. Вычисление контурных интегралов с помощью вычетов}
Пусть $a \in \mathbb{C}$ --- изолированная особая точка функции $f \in C^\omega(\mathring{U}(a), \mathbb{C})$, и пусть
$$f(z) = \sum_{n=-\infty}^{\infty} c_n (z - a)^n, \quad z \in \mathring{U}(a),$$
есть лорановское разложение $f$ в проколотой окрестности $\mathring{U}(a)$.

\textbf{Определение 11.5.} Число $c_{-1}$ называется вычетом $f$ в точке $a$ и обозначается $\text{Res}(f, a)$.

Из теоремы Лорана следует равенство
$$\text{Res}(f, a) = \frac{1}{2\pi i} \oint_{\partial D} f \, dz,$$
где $D \subset U(a)$ --- произвольная круговая окрестность $a$.

\textbf{Замечание 11.6.} Если $a$ --- устранимая особая точка, то $\text{Res}(f, a) = 0$. Если $a$ --- полюс порядка $N$, то, как следует из теоремы 11.3, функция
$$g(z) = (z - a)^N f(z)$$
голоморфна в $a$ и допускает разложение
$$g(z) = \sum_{n=0}^{\infty} a_n (z - a)^n$$
в окрестности $a$. Поскольку $c_{-1} = a_{N-1}$, имеем
$$\text{Res}(f, a) = \frac{1}{(N - 1)!} \frac{d^N}{dz^N} [(z - a)^N f(z)]. \ \blacksquare$$

\textbf{Теорема 11.7 (Коши о вычетах).} Пусть $f$ голоморфна в области $D$, за исключением конечного множества особых точек $P(f)$. Если $G \Subset D$ --- такая область, что $\partial G \cap P(f) = \varnothing$, то выполняется равенство
\begin{equation}
	\oint_{\partial G} f \, dz = 2\pi i \sum_{p \in G \cap P(f)} \text{Res}(f, p).
\end{equation}

\textbf{Доказательство.} Множество особых точек $P(f) \cap G$ конечно. Окружим каждую точку $p \in P(f) \cap G$ достаточно малой окружностью $\partial D_p$ так, чтобы все круги $\bar{D}_p$ лежали в $G$ и не пересекались друг с другом. 

Рассмотрим многосвязную область $G' = G \setminus \bigcup \bar{D}_p$. Функция $f(z)$ голоморфна в $G'$ и непрерывна на её границе. По интегральной теореме Коши для многосвязной области, интеграл по внешней границе $\partial G$ равен сумме интегралов по внутренним границам, обходимым в том же направлении (против часовой стрелки):
$$
\oint_{\partial G} f(z) \, dz = \sum_{p \in G \cap P(f)} \oint_{\partial D_p} f(z) \, dz.
$$
По определению вычета, для каждого малого контура справедливо:
$$
\oint_{\partial D_p} f(z) \, dz = 2\pi i \text{Res}(f, p).
$$
Вынося общий множитель $2\pi i$ за знак суммы, получаем искомое утверждение:
$$
\oint_{\partial G} f(z) \, dz = 2\pi i \sum_{p \in G \cap P(f)} \text{Res}(f, p). \quad \blacksquare
$$

%──────────────────────────────────────────────────────────────────────────────
\newpage
\section{Характеризация в терминах рядов Лорана изолированной особой точки $\infty$ (бесконечность). Вычет в бесконечности}

\textbf{Определение 2.} Пусть функция $f$ имеет $\infty$ своей изолированной особой точкой; ее \textit{вычетом в бесконечности} называется величина
\begin{equation}
	\mathop{\text{res}}_{\infty} f = \frac{1}{2\pi i} \int_{\gamma_R^-} f \, dz,
\end{equation}
где $\gamma_R^-$ --- окружность $\{|z| = R\}$ достаточно большого радиуса, проходимая по часовой стрелке.

Ориентация $\gamma_R^-$ выбрана так, чтобы во время ее обхода окрестность бесконечной точки $\{R < |z| < \infty\}$ оставалась слева. Напишем разложение Лорана функции $f$ в этой окрестности:
$$f(z) = \sum_{n=-\infty}^{\infty} c_n z^n.$$
Интегрируя его почленно вдоль $\gamma_R^-$ и пользуясь ортогональностью степеней, мы найдем
\begin{equation}
	\mathop{\text{res}}_{\infty} f = -c_{-1}.
\end{equation}

Члены с отрицательными степенями входят в правильную, а не в главную часть лорановского разложения в бесконечности. Поэтому, в отличие от конечных точек, вычет в бесконечности может быть не равным нулю и в том случае, когда $z = \infty$ является правильной точкой функции.

%──────────────────────────────────────────────────────────────────────────────
\newpage
\section{Логарифмический вычет, его вычисление. Приращение (полярного) аргумента вдоль пути. Принцип аргумента. Теорема Руше и ее применение}
\textbf{Определение 14.1.} Пусть функция $f$ голоморфна и локально непостоянна в проколотой окрестности $\mathring{U}(a)$ точки $a \in \mathbb{C}$. Число $\text{Res}(f'/f, a)$ называется \textit{логарифмическим вычетом} $f$ в точке $a$.

Иными словами, логарифмическим вычетом функции $f$ называется вычет её логарифмической производной $(\text{Ln} \circ f)' = f'/f$.

\textbf{Теорема 14.7 (о логарифмическом вычете).} Пусть $D \subset \mathbb{C}$ --- область, и пусть функция $f \in \mathcal{M}(D)$ локально непостоянна в $D$. Тогда для любой области $G \Subset D$ с кусочно-гладкой границей $\partial G$, не проходящей через полюсы и нули $f$, справедливо равенство
\begin{equation}
	\frac{1}{2\pi i} \oint_{\partial G} \frac{f'}{f} \, dz = N - P,
\end{equation}
где $N$ --- суммарное число нулей, а $P$ --- суммарное число полюсов функции $f$ внутри области $G$ (с учётом их кратностей и порядков).

\textbf{Доказательство.} Поскольку функция $f(z)$ мероморфна в области $D$, а область $G$ компактно вложена в $D$ ($G \Subset D$), внутри $G$ может содержаться лишь конечное число нулей и полюсов функции $f(z)$. Обозначим это конечное множество изолированных особых точек логарифмической производной через $P(f'/f) \cap G = \{a_1, a_2, \ldots, a_k\}$.

Применим основную теорему Коши о вычетах (Теорема 11.7) к функции $f'/f$ в области $G$. По условию, граница $\partial G$ не содержит нулей и полюсов, следовательно, функция $f'/f$ голоморфна на $\partial G$ и внутри $G$, за исключением точек из множества $\{a_1, a_2, \ldots, a_k\}$. Получаем:
\begin{equation}
	\frac{1}{2\pi i} \oint_{\partial G} \frac{f'}{f} \, dz = \sum_{j=1}^{k} \text{Res}\left(\frac{f'}{f}, a_j\right).
\end{equation}

Разобьем сумму в правой части равенства (2) на две группы: по всем нулям функции $f$, лежащим в $G$, и по всем её полюсам, лежащим в $G$.

1) Если точка $a_j$ является нулем функции $f$ кратности (порядка) $n_j$, то согласно Предложению 14.3:
$$ \text{Res}\left(\frac{f'}{f}, a_j\right) = n_j. $$
Суммируя по всем нулям $a_j \in G$, получаем общее число нулей с учетом кратностей:
$$ \sum_{\text{нули}} \text{Res}\left(\frac{f'}{f}, a_j\right) = \sum_{\text{нули}} n_j = N. $$

2) Если точка $a_j$ является полюсом функции $f$ порядка $m_j$, то согласно Предложению 14.3:
$$ \text{Res}\left(\frac{f'}{f}, a_j\right) = -m_j. $$
Суммируя по всем полюсам $a_j \in G$, получаем общий порядок полюсов:
$$ \sum_{\text{полюсы}} \text{Res}\left(\frac{f'}{f}, a_j\right) = \sum_{\text{полюсы}} (-m_j) = -P. $$

Подставляя эти значения обратно в равенство (2), находим:
$$ \frac{1}{2\pi i} \oint_{\partial G} \frac{f'}{f} \, dz = \sum_{\text{нули}} n_j + \sum_{\text{полюсы}} (-m_j) = N - P. \quad \blacksquare $$

Пусть $f$ --- мероморфная функция, $z : [\alpha, \beta] \to \mathbb{C}$ --- путь. Будем в зависимости от $t$ выбирать значение из $\text{Arg} \, f(z(t))$ так, чтобы функция
$$\ln f(z(t)) = \ln |f(z(t))| + i \text{Arg} \, f(z(t))$$
была непрерывна.

\textbf{Определение 14.8.} Разность
$$\Delta_z \arg f = \text{Arg} \, f(z(\beta)) - \text{Arg} \, f(z(\alpha))$$
называется \textit{приращением аргумента} функции $f$ вдоль пути $z$.

\textbf{Теорема 14.9 (принцип аргумента).} Пусть $D \subset \mathbb{C}$ --- область, и пусть функция $f \in \mathcal{M}(D)$ локально непостоянна в $D$. Тогда для любой области $G \Subset D$ с кусочно-гладкой границей $\partial G$, не проходящей через полюсы и нули $f$, справедливо равенство
$$\Delta_{\partial G} \arg f = 2\pi (N - P).$$

\textbf{Доказательство.} Действительно, полагая $\Phi = \ln \circ \, f \circ z$, имеем:
$$\oint_{\partial G} \frac{f'}{f} \, dz = \int_{\alpha}^{\beta} \Phi'(t) \, dt = \Phi(\beta) - \Phi(\alpha) = i \Delta_{\partial G} \arg f,$$
и наше утверждение следует из теоремы 14.7. $\blacksquare$

\textbf{Теорема 14.10 (Руше).} Пусть область $D \subset \mathbb{C}$ ограничена жордановой кривой, и пусть $f, g \in C^\omega(\bar{D}, \mathbb{C})$. Тогда если выполнено неравенство
\begin{equation}
	|g(z)| < |f(z)|, \quad z \in \partial D,
\end{equation}
то функции $f$ и $f + g$ имеют в $D$ одинаковое количество нулей.

\textbf{Доказательство.} Для $z \in \partial D$ справедливы неравенства
$$|f| > |g| \ge 0, \quad |f + g| \ge |f| - |g| > 0,$$
так что функции $f$ и $f + g$ не имеют нулей на $\partial D$. Но тогда
$$\Delta_\gamma \arg(f + g) = \Delta_\gamma \arg \left[ f \left( 1 + \frac{g}{f} \right) \right] = \Delta_\gamma \arg f + \Delta_\gamma \arg \left(1 + \frac{g}{f}\right) = \Delta_\gamma \arg f.$$
Применение принципа аргумента доказывает теорему. $\blacksquare$

\textbf{Следствие 14.11 (основная теорема алгебры).} Каждый многочлен $f(z) \in \mathbb{C}[z]$ степени $n$ имеет в $\mathbb{C}$ ровно $n$ корней.

\textbf{Доказательство.} Пусть
$$P(z) = \sum_{k=0}^{n} a_k z^k \in \mathbb{C}[z], \quad f(z) = a_n z^n, \quad g(z) = \sum_{k=0}^{n-1} a_k z^k,$$
так что $P = f + g$. Тогда при достаточно большом $R$ будем иметь
$$|g(z)| < |f(z)|, \quad z \in \partial U_R(0).$$
Тогда по теореме Руше
$$\sum_{z \in U_R(0)} \text{ord}(f, z) = \text{Ind}(f \circ \partial U_R(0), 0) = \text{Ind}(g \circ \partial U_R(0), 0) = n. \ \blacksquare$$

%──────────────────────────────────────────────────────────────────────────────
\newpage
\section{Теорема о среднем и принцип максимума модуля. Принцип сохранения области}

\textbf{Теорема 8.6 (теорема о среднем).} Пусть $D$ --- односвязная область в $\mathbb{C}$ и $f \in C^\omega(\bar{D}, \mathbb{C})$. Тогда
$$\forall z \in D \quad f(z) = \frac{1}{2\pi} \int_{0}^{2\pi} f(z + \rho e^{it}) \, dt,$$
если $\rho > 0$ таково, что $U_\rho(z) \subset D$.

\textbf{Доказательство.} Действительно, по интегральной формуле Коши
$$f(z) = \frac{1}{2\pi i} \oint_{\partial U_\rho(z)} \frac{f(\xi)}{\xi - z} \, d\xi = \frac{1}{2\pi i} \int_{0}^{2\pi} \frac{f(z + \rho e^{it})}{\rho e^{it}} \cdot i\rho e^{it} \, dt = \frac{1}{2\pi} \int_{0}^{2\pi} f(z + \rho e^{it}) \, dt. \ \blacksquare$$

\textbf{Теорема 14.11 (принцип максимума).} Модуль голоморфной функции не может достигать строгого локального максимума внутри области.

\textbf{Доказательство.} Пусть $f \in C^\omega(D, \mathbb{C})$ и $z_0 \in D$ --- строгий локальный максимум $|f|$ в области $D$. По теореме 8.6 о среднем,
$$|f(z_0)| = \frac{1}{2\pi} \left| \int_{0}^{2\pi} f(z_0 + \rho e^{i\varphi}) \, d\varphi \right|$$
при достаточно малом $\rho$. Но тогда
$$|f(z_0)| \le \max_{|z - z_0| = \rho} |f(z)|,$$
что противоречит определению локального максимума. $\blacksquare$

\textbf{19.6. Теорема (принцип сохранения области).} Если $f \in C^{\omega}(D,\mathbb{C})$ и $f \neq \text{const}$, то множество $f(D)$ открыто.

\textbf{Доказательство.} Возьмём произвольную точку $w_0 = f(z_0)$, $z_0 \in D$. Поскольку $D$ открыто, найдётся некоторая окрестность $B \Subset D$ точки $z_0$, при этом в силу теоремы о локальной единственности можно считать, что $f^{-1}(w_0) \cap B = \{z_0\}$.

Введём функцию $g = f - w_0$; из её непрерывности следует, что существует
\[
m \stackrel{\text{def}}{=} \min_{z \in \partial B} |g(z)| > 0.
\]
Рассмотрим $m$-окрестность $B_m \stackrel{\text{def}}{=} B(w_0, m)$ точки $w_0$ и для произвольной $w \in B_m$ введём функцию $h_w = g - w$.

Заметим, что
\[
|g(z) - h_w(z)| = |w - w_0| < m \le |g(z)|
\]
для всех $z \in \partial B$, и по теореме Руше функции $g$ и $h_w$ имеют в $B$ одинаковое число нулей, а именно 1.

Мы получили, что $g^{-1}(w) \neq \varnothing$ для любой точки $w \in B_m$, а тогда и $f^{-1}(w) \neq \varnothing$. Таким образом, точка $w_0$ лежит в $f(D)$ вместе с окрестностью $B_m$. $\quad \blacksquare$

%──────────────────────────────────────────────────────────────────────────────
\newpage
\section{Основные теоремы и приложения теории конформных отображений. Теорема Римана, принцип симметрии Римана–Шварца, принцип соответствия границ и обратный принцип соответствия границ.}

\textbf{38. Теорема Римана.} Любая односвязная область $D$, граница которой содержит более одной точки, изоморфна единичному кругу $U$.

\textbf{Теорема 1 (Риман --- Шварц).} Пусть области $D_1$ и $D_1^*$ имеют жордановы границы $\partial D_1$ и $\partial D_1^*$, причем $\partial D_1$ содержит отрезок прямой или дугу окружности $\gamma$, а $\partial D_1^*$ --- такой же отрезок или дугу $\gamma^*$. Если функция $f_1$ конформно отображает $D_1$ на $D_1^*$, причем $f_1(\gamma) = \gamma^*$, то $f_1$ допускает аналитическое продолжение в область $D_2$, симметричную с $D_1$ относительно $\gamma$, и продолженная так функция конформно отображает область $D_1 \cup \gamma \cup D_2$ на $D_1^* \cup \gamma^* \cup D_2^*$, если $D_1 \cap D_2 = \varnothing$ и $D_1^* \cap D_2^* = \varnothing$.

\textbf{Теорема (Каратеодори (Принцип соотвествия границ)).} Пусть области $D$ и $D^*$ ограничены простыми замкнутыми жордановыми кривыми $\partial D$ и $\partial D^*$. Тогда любое конформное отображение $f: D \to D^*$ можно непрерывно продолжить на границу $\partial D$ до гомеоморфизма замкнутых областей $\bar{D}$ и $\bar{D}^*$.

\textbf{Теорема (Обратный принцип соотвествия границ).} Пусть даны области $D$ и $D^*$, компактно принадлежащие $\mathbb{C}$, с жордановыми границами $\gamma$ и $\gamma^*$; если функция $f$ голоморфна в области $D$, непрерывна в $\bar{D}$ и устанавливает взаимно однозначное отображение $\gamma$ на $\gamma^*$, то и отображение $f: D \to D^*$ взаимно однозначно (т.~е. $f$ --- конформный изоморфизм).

%──────────────────────────────────────────────────────────────────────────────
\newpage
\section{Вычисление несобственных интегралов с использованием вычетов. Лемма Жордана и теорема о вычислении несобственного интеграла от рациональной функции с помощью вычетов.}

\textbf{Пример 12.5.} Рассмотрим интеграл вида
$$I = \int_{0}^{2\pi} R(\sin x, \cos x) \, dx, \quad R \in \mathbb{R}(x).$$

Сделаем в нём замену $z = e^{ix}$; тогда
\begin{equation}
	I = \oint_{|z|=1} R\left( \frac{z^2 - 1}{2iz}, \frac{z^2 + 1}{2z} \right) \frac{dz}{iz}.
\end{equation}

Если $a_1, a_2, \ldots, a_k$ --- все особые точки подынтегральной функции (12.1) в круге $|z| < 1$, то
$$I = 2\pi i \sum_{i=1}^{k} \text{Res}(f, a_i).$$

\textbf{Пример 12.6.} Пусть дана рациональная функция $f = P/Q \in \mathbb{R}(x)$, причём $Q$ не имеет действительных корней и $\deg Q \ge \deg P + 2$. Тогда
\begin{equation}
	\int_{-\infty}^{+\infty} f(x) \, dx = 2\pi i \sum_{z \in P} \text{Res}(f, z),
\end{equation}
где $P = Q^{-1}(0) \cap \{\text{Im} \, z > 0\}$.

\textbf{Доказательство.} Интеграл (12.2) существует, поскольку $f(x) = O(1/x^2)$ при $x \to \infty$. Рассмотрим полуокружность $D = \{|z| < R\} \cap \{\text{Im} \, z > 0\}$, целиком содержащую $P$. Тогда
$$\oint_{\partial D} f \, dz = 2\pi i \sum_{z \in P} \text{Res}(f, z),$$
а с другой стороны,
$$\oint_{\partial D} f \, dz = \int_{\partial D \setminus [-R, R]} f \, dz + \int_{-R}^{R} f(x) \, dx.$$

\textbf{Лемма (Жордан).} Пусть функция $f$ голоморфна в $\{\text{Im} \, z \ge 0\}$ всюду, за исключением изолированного множества особых точек, и $M(R) = \max |f(z)|$ на полуокружности $\gamma_R = \{|z| = R, \, \text{Im} \, z \ge 0\}$ стремится к нулю при $R \to \infty$ (или по последовательности $R_n \to \infty$ такой, что $\gamma_{R_n}$ не содержат особых точек $f$). Тогда для любого $\lambda > 0$
\begin{equation}
	\int_{\gamma_R} f(z) e^{i\lambda z} \, dz
\end{equation}
стремится к нулю при $R \to \infty$ (или по соответствующей последовательности $R_n \to \infty$).

(Смысл леммы состоит в том, что $M(R)$ может стремиться к нулю сколь угодно медленно, так что интеграл от $f$ по $\gamma_R$ не обязан стремиться к нулю; умножение на экспоненту $e^{i\lambda z}$ с $\lambda > 0$ убыстряет стремление к нулю.)

\textbf{Доказательство.} Обозначим через $\gamma'_R = \left\{z = R e^{i\varphi}: 0 \le \varphi \le \frac{\pi}{2}\right\}$ правую половину $\gamma_R$. В силу выпуклости синуса при $\varphi \in \left[0, \frac{\pi}{2}\right]$ имеем $\sin \varphi \ge \frac{2}{\pi} \varphi$, и, значит, на $\gamma'_R$ справедлива оценка $|e^{i\lambda z}| = e^{-\lambda R \sin \varphi} \le e^{-\frac{2\lambda R \varphi}{\pi}}$. Поэтому
$$\left| \int_{\gamma'_R} f(z) e^{i\lambda z} \, dz \right| \le M(R) \int_{0}^{\frac{\pi}{2}} e^{-\frac{2\lambda R \varphi}{\pi}} R \, d\varphi = M(R) \frac{\pi}{2\lambda} (1 - e^{-\lambda R}),$$
и интеграл по $\gamma'_R$ стремится к нулю при $R \to \infty$. Оценка для $\gamma''_R = \gamma_R \setminus \gamma'_R$ проводится аналогично. $\blacktriangleright$

%──────────────────────────────────────────────────────────────────────────────
\newpage
\section{Определение преобразования Лапласа. Теорема о существовании изображения. Поведение изображения в бесконечно удаленной точке. Изображение элементарных функций (единичная функция Хевисайда, показательная и степенная функции). Теорема обращения.}

Пусть $f : \mathbb{R} \to \mathbb{R}$, причём $f|_{\mathbb{R}_{<0}} = 0$.

\textbf{Определение 12.7.} Преобразованием Лапласа $f$ называется функция
\begin{equation}
	\mathcal{L}\{f\} : p \mapsto \int_{-\infty}^{+\infty} f(t) e^{-pt} \, dt.
\end{equation}

Её областью определения предполагается множество всех точек сходимости соответствующего интеграла. При этом функция $f$ называется оригиналом, а $F$ --- изображением. Интеграл в правой части (12.3) называется интегралом Лапласа.

\vspace{0.5em}
Обозначать связь оригинала и отображения мы часто будем так:
$$f(t) \fallingdotseq F(p), \quad \text{или} \quad F(p) \fallingdotseq f(t).$$

\textbf{Пример 12.8.} Пусть $\theta = \chi_{\mathbb{R}_{\ge0}}$ --- функция Хевисайда. Тогда $\theta \fallingdotseq 1/p$.

\vspace{0.5em}
Условимся в следующих ограничениях на оригиналы:
\begin{enumerate}
	\item $f$ кусочно-непрерывна на $\mathbb{R}$;
	\item $\exists \alpha > 0 : f(t) = O(e^{\alpha t})$ при $x \to +\infty$.
\end{enumerate}

\textbf{Предложение 12.9.} Пусть $f$ --- оригинал. Тогда функция $\mathcal{L}\{f\}$ определена и аналитична в области $\text{Re} \, z > \alpha$; соответствующий интеграл Лапласа сходится абсолютно.

\textbf{Доказательство.} Пусть $|f(t)| \le M e^{\alpha t}$ и $p = \sigma + is \in \mathbb{R} + i\mathbb{R}$. Тогда
\begin{equation}
	\int_{-\infty}^{+\infty} |f(t) e^{-pt}| \, dt \le M \int_{-\infty}^{+\infty} |e^{(\alpha - \sigma) t}| \, dt = M \left. \frac{e^{(\alpha - \sigma) t}}{\alpha - \sigma} \right|_{0}^{+\infty} = \frac{M}{\sigma - \alpha}
\end{equation}
при $\sigma > \alpha$. Из абсолютной сходимости следует и аналитичность. $\blacksquare$

\vspace{0.5em}
\textbf{Предложение 12.10.} Пусть $f$ --- оригинал. Если отображение $\mathcal{L}\{f\}$ аналитично в бесконечности, то $\lim_{p \to \infty} \mathcal{L}\{f\}(p) = 0$.

\textbf{Доказательство.} Следует из теоремы ополной сумме вычетов ввиду $\sigma = \text{Re} \, p$. $\blacksquare$

\textbf{Теорема 13.2 (теорема обращения).} Пусть $f$ --- оригинал. Тогда для любой точки непрерывности $t_0$ функции $f$ справедливо равенство
\begin{equation}
	f(t_0) = \frac{1}{2\pi i} \int_{\gamma + i\mathbb{R}} F(p) e^{pt_0} \, dp, \quad \gamma > \alpha,
\end{equation}
где $\alpha$, как и прежде, такое число, что $f(t) = O(e^{\alpha t})$ при $x \to +\infty$.


%──────────────────────────────────────────────────────────────────────────────
\newpage
\section{Основные свойства преобразования Лапласа. Теоремы линейности, подобия, затухания, запаздывания, опережения, дифференцирования и интегрирования оригинала, дифференцирования и интегрирования изображения. Свертка двух функций. Теорема умножения изображений. Доказать теоремы затухания и дифференцирования оригинала, сформулировать остальные теоремы.}

\textbf{Теорема 13.1.} Преобразование Лапласа обладает следующими свойствами.

\begin{enumerate}
	\item \textbf{Линейность.} Если $f(t) \fallingdotseq F(p)$ и $g(t) \fallingdotseq G(p)$ и $a, b \in \mathbb{C}$, то
	$$af(t) + bg(t) \fallingdotseq aF(p) + bG(p).$$
	
	\item \textbf{Подобие.} Если $f(t) \fallingdotseq F(p)$ и $\lambda > 0$, то
	$$f(\lambda t) \fallingdotseq \frac{1}{\lambda} F\left(\frac{p}{\lambda}\right).$$
	
	\item \textbf{Затухание (смещение).} Если $f(t) \fallingdotseq F(p)$ и $a \in \mathbb{C}$, то
	$$e^{at} f(t) \fallingdotseq F(p - a).$$
	
	\textbf{Доказательство.} По определению преобразования Лапласа для функции $e^{at}f(t)$ имеем:
	$$\mathcal{L}\{e^{at} f(t)\} = \int_{0}^{+\infty} \left(e^{at} f(t)\right) e^{-pt} \, dt.$$
	Используя свойства экспоненты, объединим показатели:
	$$\mathcal{L}\{e^{at} f(t)\} = \int_{0}^{+\infty} f(t) e^{-(p - a)t} \, dt.$$
	Полученный интеграл по форме абсолютно совпадает с классическим определением изображения Лапласа для функции $f(t)$, где роль комплексного параметра $p$ теперь играет выражение $(p - a)$. Следовательно:
	$$\mathcal{L}\{e^{at} f(t)\} = F(p - a). \quad \blacksquare$$
	
	\item \textbf{Запаздывание.} Если $f(t) \fallingdotseq F(p)$ и $\tau > 0$, то
	$$f(t - \tau) \fallingdotseq e^{-p\tau} F(p).$$
	
	\item \textbf{Дифференцирование оригинала.} Пусть $f$ --- кусочно-дифференцируемый оригинал. Тогда
	$$f'(t) \fallingdotseq pF(p) - f(0).$$
	
	\textbf{Доказательство.} Применим оператор Лапласа к производной $f'(t)$ по определению:
	$$\mathcal{L}\{f'(t)\} = \int_{0}^{+\infty} f'(t) e^{-pt} \, dt.$$
	Интегрируем данное выражение по частям ($\int u \, dv = uv - \int v \, du$). Пусть:
	$$u = e^{-pt} \implies du = -p e^{-pt} \, dt,$$
	$$dv = f'(t) \, dt \implies v = f(t).$$
	Подставляя эти компоненты, получаем:
	$$\mathcal{L}\{f'(t)\} = \left. f(t) e^{-pt} \right|_{0}^{+\infty} - \int_{0}^{+\infty} f(t) \left( -p e^{-pt} \right) \, dt.$$
	Рассмотрим слагаемые:
	\begin{itemize}
		\item На верхнем пределе ($t \to +\infty$) выражение $f(t) e^{-pt}$ стремится к $0$, так как для оригинала выполнено условие роста $|f(t)| \le M e^{\alpha t}$, а область определения выбирается при $\text{Re} \, p > \alpha$. На нижнем пределе ($t = 0$) получаем $f(0)e^0 = f(0)$. Значит, внеинтегральный член равен $-f(0)$.
		\item В интегральной части выносим константу $p$ за знак интеграла: 
		$$-\int_{0}^{+\infty} f(t) \left( -p e^{-pt} \right) \, dt = p \int_{0}^{+\infty} f(t) e^{-pt} \, dt = p F(p).$$
	\end{itemize}
	Собирая обе части вместе, окончательно имеем:
	$$\mathcal{L}\{f'(t)\} = pF(p) - f(0). \quad \blacksquare$$
	
	\item \textbf{Интегрирование оригинала.} Пусть $f(t) \fallingdotseq F(p)$, и пусть
	$$g(t) = \int_{0}^{t} f(\tau) \, d\tau, \quad t \in \mathbb{R}.$$
	Тогда
	$$g(t) \fallingdotseq \frac{F(p)}{p}.$$
	
	\item \textbf{Дифференцирование отображения.} Пусть $f(t) \fallingdotseq F(p)$; тогда
	$$-tf(t) \fallingdotseq F'(p).$$
	
	\item \textbf{Интегрирование отображения.} Пусть $f(t) \fallingdotseq F(p)$; тогда
	$$\frac{f(t)}{t} \fallingdotseq \int_{p}^{+\infty} F(z) \, dz,$$
	если последний интеграл сходится.
	
	\item \textbf{Теорема умножения изображений.} Пусть известны изображения $F(p)$ и $G(p)$ двух оригиналов $f(t)$ и $g(t)$. Тогда произведению этих изображений в комплексной полуплоскости соответствует свёртка исходных оригиналов в действительной области:
	$$F(p) \cdot G(p) \fallingdotseq \int_{0}^{t} f(\tau) g(t - \tau) \, d\tau.$$
	
	\item \textbf{Теорема о свёртке оригиналов.} Преобразование Лапласа от свёртки двух оригиналов $f(t)$ и $g(t)$ равно произведению их комплексных изображений:
	$$\mathcal{L}\{(f * g)(t)\} = F(p) \cdot G(p),$$
	где свёртка функций во временной области задаётся интегралом $(f * g)(t) = \int_{0}^{t} f(\tau) g(t - \tau) \, d\tau$.
\end{enumerate}

%──────────────────────────────────────────────────────────────────────────────
\newpage
\section{Три теоремы разложения. Доказать теоремы подобия и запаздывания.}

\textbf{Теорема 13.3 (первая теорема разложения).} Пусть $f$ --- оригинал, а отображение $F$ допускает разложение в ряд Лорана только по отрицательным степеням $p$:
\begin{equation}
	F(p) = \sum_{n=0}^{\infty} \frac{a_n}{p^{n+1}}, \quad |p| > R.
\end{equation}
Тогда функция $f$ раскладывается в ряд Тейлора
\begin{equation}
	f(t) = \sum_{n=0}^{\infty} \frac{a_n}{n!} t^n \quad (t > 0).
\end{equation}

\vspace{1em}

\textbf{Теорема 13.4 (вторая теорема разложения).} Если отображение $F$ является правильной рациональной дробью, то для нахождения оригинала следует разложить её на простейшие и вычислить оригинал каждого слагаемого, после чего сложить их.

\vspace{1em}

\textbf{Теорема 13.6 (третья теорема разложения).} Пусть $f$ --- оригинал, а отображение $F$:
\begin{enumerate}
	\item является аналитической функцией при $\text{Re} \, p \ge \sigma_0$,
	\item имеет при $\text{Re} \, p < \sigma_0$ лишь конечное множество особых точек $P(F)$.
\end{enumerate}
Тогда
\begin{equation}
	f(t) = \sum_{q \in P(F)} \text{Res}\Big( p \mapsto e^{pt}F(p), \, q \Big).
\end{equation}

\textbf{Подобие.} Если $f(t) \fallingdotseq F(p)$ и $\lambda > 0$, то
$$f(\lambda t) \fallingdotseq \frac{1}{\lambda} F\left(\frac{p}{\lambda}\right).$$

\textbf{Доказательство.} По определению преобразования Лапласа для функции $f(\lambda t)$ имеем:
$$\mathcal{L}\{f(\lambda t)\} = \int_{0}^{+\infty} f(\lambda t) e^{-pt} \, dt.$$
Сделать замену переменной интегрирования. Пусть $u = \lambda t$, тогда $t = \frac{u}{\lambda}$ и $dt = \frac{1}{\lambda} du$. Поскольку $\lambda > 0$, пределы интегрирования остаются прежними (от $0$ до $+\infty$). Подставим новую переменную в интеграл:
$$\mathcal{L}\{f(\lambda t)\} = \int_{0}^{+\infty} f(u) e^{-p \left(\frac{u}{\lambda}\right)} \frac{1}{\lambda} \, du.$$
Вынесем константу $\frac{1}{\lambda}$ за знак интеграла и сгруппируем переменные в показатели экспоненты:
$$\mathcal{L}\{f(\lambda t)\} = \frac{1}{\lambda} \int_{0}^{+\infty} f(u) e^{-\left(\frac{p}{\lambda}\right) u} \, du.$$
Полученный интеграл по форме абсолютно совпадает с определением изображения функции $f$, где роль комплексного параметра $p$ теперь играет величина $\frac{p}{\lambda}$. Следовательно:
$$\mathcal{L}\{f(\lambda t)\} = \frac{1}{\lambda} F\left(\frac{p}{\lambda}\right). \quad \blacksquare$$

\textbf{Запаздывание.} Если $f(t) \fallingdotseq F(p)$ и $\tau > 0$, то
$$f(t - \tau) \fallingdotseq e^{-p\tau} F(p).$$

\textbf{Доказательство.} Применим оператор Лапласа к запаздывающей функции $f(t - \tau)$ по определению:
$$\mathcal{L}\{f(t - \tau)\} = \int_{0}^{+\infty} f(t - \tau) e^{-pt} \, dt.$$
Поскольку по определению оригинала $f(t - \tau) = 0$ при $t < \tau$, интеграл на промежутке $[0, \tau)$ равен нулю. Изменим нижний предел интегрирования на $\tau$:
$$\mathcal{L}\{f(t - \tau)\} = \int_{\tau}^{+\infty} f(t - \tau) e^{-pt} \, dt.$$
Сделать замену переменной сдвига времени $u = t - \tau$, откуда $t = u + \tau$ и $dt = du$. При этом нижний предел $t = \tau$ переходит в $u = 0$, а верхний предел остается бесконечным. Подставим новые переменные:
$$\mathcal{L}\{f(t - \tau)\} = \int_{0}^{+\infty} f(u) e^{-p(u + \tau)} \, du.$$
Распишем экспоненту как $e^{-p(u + \tau)} = e^{-pu} \cdot e^{-p\tau}$. Поскольку множитель $e^{-p\tau}$ не зависит от переменной интегрирования $u$, вынесем его за знак интеграла:
$$\mathcal{L}\{f(t - \tau)\} = e^{-p\tau} \int_{0}^{+\infty} f(u) e^{-pu} \, du.$$
Оставшийся интеграл представляет собой классическое определение изображения $F(p)$ для оригинала $f(t)$. Таким образом, получаем:
$$\mathcal{L}\{f(t - \tau)\} = e^{-p\tau} F(p). \quad \blacksquare$$
\end{document}

